\documentclass[11pt,onecolumn]{article}
\usepackage[margin=1.5in]{geometry}
\usepackage{times}
\usepackage{setspace}
\setstretch{1.5}
\usepackage{amsmath,amssymb}
\usepackage{array}
\usepackage{booktabs}
\usepackage[hidelinks]{hyperref}
\usepackage{microtype}

\usepackage{fancyhdr}
\pagestyle{fancy}
\fancyhf{}
\fancyhead[L]{\footnotesize Review copy}
\fancyhead[R]{\footnotesize Breaking the HBM wall}
\fancyfoot[C]{\footnotesize\thepage}

\setlength{\parskip}{0.3em}
\setlength{\parindent}{0em}

\title{Breaking the HBM wall: A Distributed Spatial Processing-Near-Memory Architecture using DUV ASICs and Deterministic Routing}
\author{Bowen Gu}
\date{August 23, 2026 (review copy)}

\begin{document}
\maketitle
\thispagestyle{empty}

\section*{Abstract}\label{sec:abstract}

The memory wall has hardened into a physical wall: High-Bandwidth Memory capacity per package is bounded by interposer reticle limits\cite{15}, and the manufacturing economics behind HBM stacks---which consume roughly three times the wafer area per bit of commodity DRAM, at extreme power density---are structural rather than cyclical\cite{25}, just as the workload mix shifts toward capacity-bound jobs---including Mixture-of-Experts transformer inference, scientific stencil computation, climate modeling, seismic imaging, and large-scale numerical simulation---whose binding resource is resident memory, not peak FLOPs\cite{3,4,26}. This paper proposes a distributed spatial Processing-Near-Memory (PNM) architecture that bypasses the HBM wall with commodity 192-bit LPDDR6 CAMM2 modules\cite{16,17} paired with mature-node DUV MAC ASICs, wired by a deterministic single-spine wormhole fabric of stateless Hardware Flit Repeaters and combinational lxy\_repeater nodes that replaces runtime scheduling and cache coherency with coordinate arithmetic and closed-form latency bounds\cite{13,14}. A 512-node reference chassis provides 64 TB of attached memory, 131 TB/s aggregate node-local bandwidth, and approximately 197 TFLOPS FP64 at approximately 15 W per node---8--10 kW per chassis against 0.25 MW for capacity-equivalent HBM---managed by one socketed RISC-V BMC/router chip handling topology discovery, workload dispatch, and host PCIe\cite{5}. The result is bit-exact-reproducible computation with provable deadlock freedom, compile-time latency bounds, and a fraction of the thermal envelope, interposer dependency, and manufacturing complexity of monolithic GPU accelerators. Validation scope is stated explicitly: the machine-checked claims are byte-exact delivery and replay, doorbell refusal accounting, and cycle-exact closed-form latency on a 100 MHz behavioral co-simulation (Sections \hyperref[sec:3.5]{3.5} and \hyperref[sec:4.4]{4.4}); bandwidth, power, and nanosecond latencies are projections of characterized link and process assumptions, not silicon measurements.

Keywords: Processing-Near-Memory; Wormhole Switching; Deterministic Routing; Dimension-Order Routing; Spatial Computing; LPDDR6 CAMM2; DUV ASIC; Mixture-of-Experts Transformers; Stencil Computation; Roofline Analysis; RISC-V; BMC; High-Performance Computing

\section*{Introduction}\label{sec:introduction}

\subsection*{1.1 The physical bottleneck}\label{sec:1.1}

The prevailing paradigm for large-scale transformer inference relies on monolithic GPU accelerators tightly coupled to High-Bandwidth Memory (HBM) via advanced silicon interposers\cite{2,15}. This approach is hitting three hard physical limits that no process shrink or market cycle can resolve. First, HBM capacity per package is bounded by interposer reticle limits: each HBM stack is a vertical assembly of DRAM dies connected through silicon vias, and the interposer that bridges these stacks to the GPU die has a fixed reticle-bounded area\cite{15}. This caps the number of stacks that can be placed around a single GPU, which in turn caps the aggregate memory capacity per accelerator---currently 80 GB for the largest shipping parts. For capacity-bound workloads this ceiling is the binding constraint: the compute exists to process the data set, but the memory to hold it does not fit. Mixture-of-Experts (MoE) transformers are the most visible instance---hundreds of billions of parameters, most of them idle during any given forward pass\cite{3,4}---but the same profile recurs across sparse retrieval, scientific simulation state, and seismic image volumes. Second, the manufacturing burden of HBM is concentrated and severe: HBM stacks consume roughly three times the wafer area per bit of commodity DRAM, and the 2025--26 DRAM shortage---driven by fabrication capacity shifting from commodity DRAM toward HBM---demonstrates the supply fragility this creates\cite{25}. Third, the power density of an advanced interposer stack is extreme: an H100-class GPU dissipates roughly 300 W in a package smaller than a paperback book, and the thermal management of stacked DRAM dies---which cannot tolerate junction temperatures above 85$^\circ$C---limits the number of layers that can be reliably cooled. A single NVIDIA H100-class GPU with 80 GB of HBM can cost upwards of thirty thousand dollars---roughly \$375 per gigabyte of attached memory---yet the majority of its transistor budget and thermal headroom is spent not on computation, but on moving data across hierarchical memory caches and network interfaces\cite{2}. The architecture described in this paper sidesteps all three limits by replacing interposer-bound HBM with commodity LPDDR6 CAMM2 modules\cite{16,17}---each a standard JEDEC part with no interposer dependency, no advanced lithography requirement, and a thermal envelope of approximately 5 W---paired with a mature-node 14 nm or 28 nm DUV MAC ASIC. The Results section quantifies the alternative: a 512-node chassis provides 64 TB at approximately 15 W per node and approximately 8--10 kW per chassis, versus 0.25 MW and 819 GPUs for capacity-equivalent HBM.

\subsection*{1.2 Design philosophy: hardware as a static dataflow graph}\label{sec:1.2}

This paper advances an alternative philosophy inspired by Unix pipelines\cite{7}, pure functional programming\cite{6}, operating-system microkernels\cite{8}, and the dataflow computing tradition\cite{11,12}. We treat the physical motherboard not as a passive substrate for a Von Neumann processor, but as an active, immutable dataflow graph. Each Hardware Flit Repeater (HFR) does exactly one thing---it forwards a flit to the next node---just as a Unix program does one thing well and composes via pipes. The hardware itself enforces functional immutability: flits are immutable once injected, HFRs are stateless functions of their input, and MAC ASICs produce deterministic outputs from static memory contents. There are no runtime page tables, no dynamic branch prediction, no cache coherency protocols, and no global mutable state anywhere in the system. The entire machine is a single, spatially distributed pure function---a physical instantiation of a dataflow graph---making it a natural substrate for first-order, statically allocated functional programs (Section \hyperref[sec:2.14]{2.14}). The lxy\_repeater extends the microkernel principle from software to silicon: it enforces a strict, non-bypassable physical separation between layers through a hardware bitmask comparison, with no configuration register to misprogram, no firmware to exploit, and no driver to crash\cite{9}. This purity is not merely aesthetic: it is what makes bit-exact replay after faults (Section \hyperref[sec:4.1]{4.1}) and compile-time latency bounds (Section \hyperref[sec:4.3]{4.3}) possible at all.

\subsection*{1.3 Manufacturing realities: DUV versus EUV}\label{sec:1.3}

A common objection to custom accelerator silicon is the perceived necessity of bleeding-edge process nodes. This architecture rejects that premise. Memory bandwidth, not transistor switching speed, is the binding constraint for inference and stencil workloads. The roofline analysis of Section \hyperref[sec:3.2]{3.2} shows both flagship workload families sit deep in the bandwidth-bound regime at the reference machine balance\cite{5}, so a 14 nm MAC array never limits performance. A 14 nm or 28 nm Deep Ultraviolet (DUV) ASIC is therefore perfectly sufficient to saturate a local LPDDR6 bus\cite{17}, and mature nodes offer structural advantages that no market cycle can erase: dramatically higher yields than EUV lithography, established supply chains with predictable lead times, and the freedom to use large die areas for MAC arrays without the exponential cost curve of sub-7 nm fabrication. The 2025--26 DRAM shortage---driven by fabrication capacity shifting from commodity DRAM toward HBM, which consumes roughly three times the wafer area per bit\cite{25}---further illustrates the fragility of an ecosystem built on interposer-bound, EUV-dependent memory. By pairing commodity LPDDR6 CAMM2 modules with mature-node logic, the architecture decouples memory capacity from the interposer and lithography bottlenecks that constrain HBM.

\subsection*{1.4 ASIC philosophy and historical precedents}\label{sec:1.4}

General-purpose processors derive flexibility from massive instruction decoders, register files, branch predictors, and cache hierarchies. An ASIC burns transistors only on the logic that directly contributes to the target computation\cite{10}. Our architecture applies this principle to mature DUV nodes: the DUV MAC ASIC contains no instruction decoder, no branch predictor, and no operating-system interface---only a fixed-function MAC array hard-wired to its local LPDDR6 bus. The design inherits the INMOS Transputer's philosophy of computing elements with direct, deterministic communication\cite{12}, but replaces its general-purpose processors with fixed-function MAC ASICs and its serial links with high-bandwidth LPDDR6 buses. The wormhole routing fabric draws on Dally and Seitz's work on deadlock-free routing: deterministic routing is deadlock-free if and only if the channel dependency graph is acyclic\cite{13,14}. Our fabric guarantees acyclicity by construction through monotonically ordered virtual-channel classes (Section \hyperref[sec:4.4]{4.4}).

\section*{Methods}\label{sec:methods}

\subsection*{2.1 The Z-axis spine and repeater nodes}\label{sec:2.1}

The network topology is a strictly deterministic single-spine tree---a spine-and-leaf fabric with exactly one spine. A central Z-axis spine runs vertically through the chassis, bridging parallel horizontal X/Y motherboards with high-density mezzanine connectors. Traffic between nodes on the same layer follows exactly one dimension-ordered X/Y path; traffic between layers additionally enters the spine at its own layer's attachment, travels monotonically along the spine to the destination layer's attachment, and continues across the destination board's dimension-ordered path. Because there is exactly one legal route between any pair of nodes, routing requires no search algorithm, no arbitration, and no adaptive logic---only coordinate arithmetic.

At each layer, an lxy\_repeater evaluates the incoming flit header against a local hardware ID bitmask. Implemented as a combinational tree of XOR and AND gates on a 14 nm DUV node, this comparator operates in approximately 50--100 picoseconds without clocks, queues, or software intervention in the compare path. The comparator's match decision is registered on the local Network-on-Board (NoB) clock edge; clock-domain crossings use standard source-synchronous forwarding between HFRs. Matching flits are gated onto the local X/Y NoB; non-matching flits proceed along the spine. The lxy\_repeater boundary acts as a rigid physical microkernel barrier\cite{9}. Direction is decided ahead of time, never at runtime: the AOT compiler (Section \hyperref[sec:2.8]{2.8}) emits a deterministic, immutable routing table for every repeater, loaded into the repeaters beside each ASIC node during initialization; each lxy\_repeater steers flits solely from its pre-loaded table, with nothing to adapt, arbitrate, or reconfigure on the data path.

The spine is the fabric's only bisection, and it is sized by arithmetic, not by assertion. Let $\alpha$ be the cross-layer traffic fraction of the dominant workload---a compile-time constant derived by the AOT mapper from the resident placement of experts or stencil subdomains (Section \hyperref[sec:2.8]{2.8}). If the dominant cross-layer trip moves $S$ bytes (a kilobyte-class token plus its returned hidden state, $S \approx 8$ KB), then a spine provisioned at $B$ bytes per second per direction sustains at most $R = B / (\alpha S)$ cross-layer dispatches per second. The reference chassis provisions the spine as $\approx$128 shared parallel 16 GB/s lanes ($\approx$2 TB/s aggregate per direction, shared by all 512 nodes), which sustains $\approx$3$\times10^{8}$ cross-layer dispatches/s under naive random expert placement ($\alpha = 7/8$) and rises past $10^{9}$ dispatches/s once hot-expert replication and layer-affine placement cut $\alpha$ to $\approx$1/4 (Section \hyperref[sec:2.12]{2.12}). Each generated token costs 270 fabric dispatches---30 dense attention projections plus 240 MoE expert activations for the Gemma-4-class reference workload (Section \hyperref[sec:2.6]{2.6})---so the honest end-to-end token rates are $\approx$1.1$\times10^{6}$ tokens/s at $\alpha = 7/8$ and $\approx$3.7$\times10^{6}$ tokens/s at $\alpha \approx 1/4$: the dispatches/s figures are fabric throughput, and the token rate follows by dividing by the per-token dispatch count. The rule is exact because every route is known to the compiler. Intra-layer stencil halo traffic never enters the spine (Section \hyperref[sec:2.2]{2.2}). A fat-spine variant---adding lanes toward the root---scales this bound for larger stack heights at the cost of connector pins, which are the dominant spine cost.

The system therefore reduces to four kinds of hardware. Commodity DRAM (192-bit LPDDR6 CAMM2) holds all weights and state. A DUV MAC ASIC beside each module performs the only computation. Stateless Hardware Flit Repeaters (HFRs) and the combinational lxy\_repeaters form the entire data path. And one central router chip---the only programmable silicon in the machine---owns everything that is not streaming data: it performs Power-On Self-Test and topology discovery (Section \hyperref[sec:2.5]{2.5}), terminates the host PCIe uplinks (Section \hyperref[sec:4.2]{4.2}), evaluates the coordinate mapping for dynamic MoE expert routing (Section \hyperref[sec:2.8]{2.8}), and injects every JIT-dispatched token into the fabric. Everything on the hot path is pure transport; everything with state lives in a single chip at the spine root.

\subsection*{2.2 The intra-layer fabric: dimension-order routing}\label{sec:2.2}

Within a layer, HFRs form a direct X/Y grid that mirrors the motherboard's physical socket layout. On-board routes use X-then-Y dimension order: a flit travels along the X axis to the destination column, then along the Y axis to the destination socket\cite{13,14}. Dimension-order routes are deterministic, minimal, and free of cyclic channel dependencies by the standard results. A halo exchange between neighboring grid cells is a single-hop route between adjacent sockets that never touches a lxy\_repeater or the spine.

Every fabric link---on-board and spine alike---is characterized as a 128-bit-wide source-synchronous channel at 1 GHz, providing $\approx$16 GB/s per link. The HDL sketch and the co-simulation harness model the byte-oriented flit stream at one byte per cycle on a 100 MHz testbench clock; the 1 GHz 128-bit physical link is never clocked in simulation. The byte-serial 100 MHz model is therefore slower than the characterized link by roughly 160$\times$---a 10$\times$ clock gap times sixteen bytes per cycle---understating per-link bandwidth by that factor while preserving the routing-decision logic that is the point of the model. The nanosecond latencies of Section \hyperref[sec:3.4]{3.4} are projections of the closed-form hop-count expression at the characterized 16 GB/s link rate, not measurements of the 100 MHz testbenches; the testbenches verify routing, framing, and the doorbell discipline cycle-exactly at their own clock, and Table 3 reports those cycle counts as the verification record. Cross-node traffic is bounded by these 16 GB/s links, which is deliberately separate from the $\approx$256 GB/s each node can read from its own memory: the aggregate figure in Section \hyperref[sec:3.1]{3.1} is a sum of node-local bandwidths, and any workload that must move data between nodes is fabric-bound at the link rate.

\subsection*{2.3 The motherboard and isothermal manifold}\label{sec:2.3}

The horizontal plane consists of rigid X/Y motherboards populated with paired LPDDR6 CAMM2 memory sockets and DUV MAC ASIC sockets (Section \hyperref[sec:2.4]{2.4}). Stacked vertically, these boards create a severe thermal challenge: LPDDR6 modules throttle above 85$^\circ$C. At the reference operating point, each node dissipates approximately 15 W (about 5 W DRAM, 9 W MAC array, 1 W fabric share), a fully populated 64-node board about 1 kW, and an eight-layer chassis about 8--10 kW including fabric and control losses. The chassis incorporates an Isothermal Parallel Manifold---rigid vertical fluid channels that inject identical-temperature coolant (30$^\circ$C) simultaneously across all stacked layers. Balance is enforced, not assumed: each branch is fitted with a calibrated flow orifice sized during chassis assembly, holding each branch's coolant temperature rise within a $\pm$2$^\circ$C band. At approximately 15 L/min per chassis (water, 10$^\circ$C coolant temperature rise), the worst-case DRAM junction temperature remains at or below 70$^\circ$C---15$^\circ$C below the throttling threshold.

\subsection*{2.4 CAMM2 integration, packaging, and sideband signaling}\label{sec:2.4}

Each compute node centers on a JEDEC LPDDR6 CAMM2 memory module (the socketable form factor that succeeds JESD318's LPCAMM2)\cite{16,17}. LPDDR6 (JESD209-6) organizes its bus as 24-bit channels, each made of two 12-bit sub-channels; the CAMM2 module is 8$\times$24 bits (four 48-bit channel groups), i.e. a **192-bit** interface. The reference configuration assumes this 192-bit bus at 10.7 Gb/s/pin, yielding $\approx$256 GB/s per node. The DUV MAC ASIC sits adjacent to---rather than inside---the DRAM die, making this Processing-Near-Memory (PNM) rather than true Processing-in-Memory (PIM). A JEDEC LPDDR6 interface requires a host-side memory controller, and this design integrates that controller into the DUV MAC ASIC itself: the controller lives on the motherboard in the ASIC socket beside each CAMM2 module, its LPDDR6 PHY pins of Table 1 being the memory side of the interface (Section \hyperref[sec:2.16]{2.16}), with no separate controller chip on the board. Two packaging variants define adjacency:

Variant A --- Socketed ASIC on the motherboard (reference design). Each CAMM2 socket is paired with an adjacent low-profile ASIC socket (land-grid-array or card-edge mezzanine). The ASIC is socketed rather than soldered: nodes can be populated, depopulated, and replaced in the field, matching the POST discovery sequence (Section \hyperref[sec:2.5]{2.5}) and the failure model (Section \hyperref[sec:4.1]{4.1}). The memory module remains an unmodified commodity part. Placing the two sockets within roughly 20 mm keeps the LPDDR6 interface short enough for commodity-grade signal integrity.

Variant B --- Module-integrated ASIC (high-density option). The DUV MAC ASIC is mounted directly on the CAMM module substrate beside the DRAM packages. This yields the shortest memory interface and lowest interface power, at three costs: the module becomes a custom part, ASIC heat sits adjacent to the DRAM, and a failed ASIC retires the entire module.

Direct BGA attachment of the ASIC to the motherboard was considered and rejected: a permanent solder joint defeats field serviceability and contradicts the reconfigurable-inventory philosophy of Section \hyperref[sec:2.5]{2.5}. The reference configuration in the Results section assumes Variant A.

Because the CAMM2 connector pins are fully occupied by the memory bus, four sideband signals travel on a separate low-pin-count board-level sideband header (in Variant B, a short flex tail from the module substrate). Each is a single wire from the node, and each runs as part of a daisy-chained control net that climbs the spine to the central router chip; the net is a low-speed control plane entirely separate from the data-path NoB:

DOORBELL\_TRIG: asserted by the node's DMA engine only when (a) the received byte count equals the message-length field in the flit header, (b) the end-to-end message CRC validates, and (c) the CRC-protected destination field equals the node's own coordinate. Because HFR forwarding along a single path is strictly in-order, byte-count equality detects complete delivery; a stream aborted mid-message can never satisfy all three conditions. The full state machine is specified in Section \hyperref[sec:2.9]{2.9}.

DOORBELL\_ACK: asserted by the same DMA engine in the same cycle as DOORBELL\_TRIG, and forwarded up the sideband net. It is the source's completion signal, and its single-wire cost is the entire acknowledgment overhead of the design: there is no ACK packet, no sequence number, and no extra latency on the data path. A source that has not seen DOORBELL\_ACK within its compile-time latency budget infers a lost or un-executable message and raises NODE\_ERR for that destination (Section \hyperref[sec:4.1]{4.1}). The ACK wire is the loss detector that turns the doorbell's refusal-to-fire into a recoverable, observable event.

NODE\_ERR: a fatal fault line asserted on the sideband control net, streaming directly to the central router chip (Section \hyperref[sec:4.1]{4.1}). It does not travel on the data path, so it remains usable when the node's NoB or DMA path is dead.

TOPOLOGY\_RDY: asserted during Power-On Self-Test once voltage margins and link training succeed.

Board outline and pin budgets. The motherboard's wiring is entirely point-to-point and length-deterministic---the property that makes the closed-form latencies of Section \hyperref[sec:3.4]{3.4} possible. Table 1 itemizes the four connector populations a node touches; the counts are pin-budget outlines, not JEDEC exactness, since the LPDDR6 CAMM2 class is forward-looking while the current LPCAMM2 class (JESD318) is a 644-pin module\cite{16}. The ASIC side of the same budget, with full electrical specifications, is tabulated in Section \hyperref[sec:3.6]{3.6}.

Table 1. Board-level pin and trace outline (per node; first-order estimates).

Interface | Pins | Routing
CAMM2 socket (LPDDR6, 192-bit) | $\approx$600--650 | 192 DQ, 24 DQS pairs, CA/CS/CK, power/ground; $\leq$20 mm to the ASIC socket
ASIC socket (Variant A) | $\approx$800--850 | LPDDR6 PHY $\approx$290, fabric 128-bit $\times$2 $\approx$270, sideband 4, clock/test $\approx$20, power/ground remainder
Sideband header | 6 | 4 signal + ground (+3.3 V); daisy-chained control net up the spine
Spine mezzanine | 128 $\times$ lanes + control | per-link 128-bit source-synchronous with forwarded clock

Every trace is impedance-controlled and length-matched within its group. The LPDDR6 escape from the CAMM2 socket to the ASIC socket runs $\leq$20 mm on dedicated DDR routing layers as a tightly matched bundle---the shortest legal memory route on the board. The fabric links are point-to-point HFR-to-HFR runs on their own layers, driven source-synchronously: the clock is forwarded beside the data, so per-link skew is absorbed by the receive FIFO and the wire delay is the single deterministic term the compiler's closed form uses. A 64-node board is estimated at 10--12 routing layers: two DDR escape planes, two fabric planes, a spine-breakout region, and power/ground planes between them. The four sideband signals share one low-speed daisy-chain trace pair climbing the spine to the central router chip; their single-wire-per-signal cost is the entire control-plane wiring budget of the design.

\subsection*{2.5 Initialization and the POST Discovery Sequence}\label{sec:2.5}

The firmware state machine boots the chassis through five deterministic phases, each completing before the next begins. The central router chip drives the sequence; no node participates until the phase that addresses it. The state machine is:

RESET $\rightarrow$ POST\_DISCOVERY $\rightarrow$ ROUTING\_TABLE $\rightarrow$ WEIGHT\_UPLOAD $\rightarrow$ MOE\_LOAD $\rightarrow$ READY

Phase 1: POST\_DISCOVERY. The router emits a ping frame; each layer's lxy\_repeater propagates it along the local X/Y NoB. Every populated node that completes link training and Built-In Self-Test asserts TOPOLOGY\_RDY. The router aggregates responses into a live inventory: for each coordinate $(L, X, Y)$, the inventory records the 8-bit Module ID ($X[3:0] \parallel Y[3:0]$), the per-link bandwidth ($\approx$256 GB/s from LPDDR6 CAMM2), and the node's status (ready or down). At the reference scale, 512 LPDDR6-class links train in milliseconds each and run in parallel; the ping/inventory sweep completes in microseconds, and the router aggregates the inventory in well under one second before any weight is loaded. The inventory is the Instruction Set Architecture for this boot instance: it defines the set of live coordinates, the physical topology, and the deterministic latency between every pair.

Phase 2: ROUTING\_TABLE. The router programs every lxy\_repeater and HFR in the chassis with its routing-table entry. Each entry is an 11-bit routing bitmap: bits [10:7] hold the 4-bit Layer ID (1-based), bit [6] selects the AXIS (0 = X, 1 = Y), bit [5] selects the SIGN (0 = +, 1 = $-$), and bits [4:0] hold the hop distance from the lxy\_repeater. The router computes these bitmaps from the inventory: for a node at $(L, X, Y)$, the lxy\_repeater at layer $L$ receives a bitmap matching Layer $L$; the node itself receives a bitmap matching its full coordinate with AXIS = Y and DIST = $Y$. The bitmaps are loaded via the sideband control net during this phase; the fabric's deterministic single-path structure means every bitmap is a function of the two coordinate pairs, not a search result. In the co-simulation harness, the bitmaps are embedded in the generated Verilog topology (Section \hyperref[sec:2.15]{2.15}), so this phase is a no-op; in production, the sideband daisy chain carries the bitmaps from the router to each repeater and node.

Phase 3: WEIGHT\_UPLOAD. The router streams weight tensors to every node's LPDDR6 bank. For each node, the router constructs a weight-upload command: a flit containing the target Layer ID, Module ID, payload length, and the weight bytes (synthesized from the model compiler's output, Section \hyperref[sec:2.7]{2.7}). The wire format is:

CMD (0x01) | LAYER | MODULE\_ID | LEN\_LO | LEN\_HI | payload... | CRC\_HI | CRC\_LO

The payload carries the weight data for the node's assigned compute unit: BF16 FMA weights for MoE experts, BF16 systolic array weights for attention projections, FP32 constants for layer norms, and INT8 weights for quantized inference. Each upload command also records the compute unit type (CUType), the data type (DType), the model layer index, and the tensor role (expert, attention, layernorm, embedding) so the host-side driver can verify weight upload integrity per unit type. The router sorts commands by $(L, X, Y)$ for deterministic ordering and streams them in lexicographic coordinate order, ensuring that nodes on the same layer receive their weights before the next layer's upload begins. A fully populated 64 TB chassis loads in under ten minutes at uplink line rate; the reference design also supports per-node boot flash on the sideband header (roughly two minutes through 512 parallel streams).

Phase 4: MOE\_LOAD. The router loads the MoE gating state into its own SRAM. Expert selection lives exclusively on the router chip, which evaluates it at runtime to determine which expert handles each token. The RTL's moe\_gating unit (Section \hyperref[sec:2.17]{2.17}) implements the learned-projection datapath---logits $=$ hidden $\times$ router\_weights followed by top-$k$ selection---with projector weights streamed over PCIe into its weight SRAM. Because this work ships no end-to-end-trained projector, the dispatch path the co-simulation actually exercises routes each token through a deterministic token-seeded routing hash evaluated in router firmware (Section \hyperref[sec:2.12]{2.12}), which produces every routing decision quoted in Section \hyperref[sec:3.5]{3.5}; the RTL projection datapath and the hash share the top-$k$ selection shape but not the learned behavior or expert identity. A learned $128$-expert BF16 projector at the reference model's gating dimension needs roughly $0.3$--1 MB of weights---5--15$\times$ the router's dedicated 64 KB gating SRAM (32K BF16 entries)---so learned-router sizing and its per-token evaluation latency are projections (Section \hyperref[sec:4.4]{4.4}). In the co-simulation, this phase records the gating state as loaded; in production, the weights arrive over PCIe from the host before execution begins.

Phase 5: READY. The chassis is fully initialized. The router exports a topology schema describing physical coordinates, per-link bandwidth, and deterministic latency between each pair---the closed-form latency is hop count $\times$ per-hop delay plus message serialization at the link rate (Section \hyperref[sec:3.4]{3.4}), not a search result. The AOT compiler (Section \hyperref[sec:2.8]{2.8}) then lowers the user IR against this exact schema, so the ISA for this boot instance exists within a second of power-on. The router enters the dispatch loop: for MoE workloads, it evaluates the gating network and dispatches tokens; for stencil workloads, it triggers halo exchanges at the scheduled timestamps. The dispatch loop processes each token through every model layer in sequence: for each layer, the firmware first dispatches to the dense attention node (round-robin by model layer), then checks whether any KV cache bank on that layer exceeds the offload threshold---if so, it dispatches a KV offload command and triggers a reload cycle---and finally dispatches to the top-$k$ MoE expert nodes identified by the gating network. Each dispatch constructs a wormhole flit carrying the [Layer ID | Module ID] header, a CTRL byte encoding the virtual-channel class, the token payload, and the end-to-end CRC, and injects it into the fabric. The dispatch record tracks the layer, phase (dense, moe, or kv\_offload), target node, expert index, flit byte count, KV action, compute unit type, and tensor role for host-side verification.

The entire boot sequence is deterministic and bounded. Every phase completes in known time: discovery in $<$1 s, routing-table load in $<$1 s (512 bitmaps at 6 bytes each via sideband), weight upload in $<$10 min (64 TB at uplink line rate), MoE load in $<$1 ms. The chassis does not enter the READY state until every phase completes successfully; if any phase fails (link training timeout, CRC mismatch on a weight upload, routing-table load NAK), the router raises NODE\_ERR for the affected coordinate and continues in degraded mode (Section \hyperref[sec:4.1]{4.1}).

\subsection*{2.6 Software Stack and Execution Model}\label{sec:2.6}

The software stack is a narrow, layered pipeline that compiles high-level source languages onto the physical chassis. Unlike conventional GPU software stacks---which involve runtime drivers, kernel-mode scheduling, memory management units, and dynamic allocation---the PNM stack is entirely ahead-of-time: every weight, every route, and every dispatch timestamp is fixed before execution begins. The stack has four layers.

Layer 1: Source-language front ends. The architecture supports three source-language compilers that emit a typed intermediate representation (IR) targeting the chassis's compute units. The R compiler (CompileR) parses a subset of R---assignments, arithmetic, aggregations (sum, mean, min, max), comparisons, and function calls---and emits FP64 IR instructions targeting the `fp64\_fma` unit. The Haskell compiler (CompileHaskell) parses function definitions, guards, pattern matching, and higher-order aggregations (sum, product, minimum, maximum) and emits the same FP64 IR, supporting recursive function inlining and user-defined function composition. The HLSL compiler (CompileHLSL) parses GPU shading language---float declarations, arithmetic, intrinsic calls (dot, lerp, clamp, saturate, smoothstep, step, rcp, sqrt), and vector constructors---and emits FP32 ALU IR instructions targeting the `fp32\_alu` unit. All three compilers are deterministic, side-effect-free, and produce IR that the downstream pipeline consumes without modification.

Layer 2: The intermediate representation. The IR is a typed, register-based instruction set with two precision domains. The FP64 domain (R and Haskell) provides: f64.load, f64.store, f64.add, f64.sub, f64.mul, f64.fma, f64.neg, f64.div, f64.min, f64.max, f64.cmp, f64.mov, and f64.const. The FP32 ALU domain (HLSL) provides: alu.load, alu.store, alu.const, alu.add, alu.sub, alu.mul, alu.div, alu.min, alu.max, alu.cmp, alu.dot, alu.lerp, alu.clamp, alu.rcp, and alu.mov. The IR is deliberately minimal: no loops, no branches, no dynamic memory allocation. This constraint is what makes the downstream AOT pipeline sound: every instruction maps to a fixed-function hardware unit, and the compiler can reason about the entire execution trace at compile time. Complex control flow---loops, conditionals, recursion---is lowered by the source-language front end into straight-line IR with explicit data dependencies, which the hardware executes as a static dataflow graph.

Layer 3: The AOT compilation pipeline (Section \hyperref[sec:2.7]{2.7}). The five-stage pipeline---ingestion, partitioning, kernel mapping, coordinate assignment, and dispatch scheduling---takes the IR and the model weights and produces a chassis schema and a `.pnm` program. The compiler assigns each IR instruction to a specific node's compute unit, routes the data dependencies through the fabric, and schedules every message injection to respect buffer credits and link bandwidth. The output is a static, immutable execution plan that the hardware follows without runtime indirection.

Layer 4: The firmware and runtime. The central router chip's firmware (Section \hyperref[sec:2.5]{2.5}) boots the chassis, loads weights, and runs the READY-phase dispatch loop, managing the KV cache (Section \hyperref[sec:2.12]{2.12}) with least-recently-used eviction when capacity is exhausted. A C port targets MCU implementations (ARM Cortex-M/R, RISC-V) for embedded deployments, using static allocation only---no `malloc`, no dynamic memory, no operating-system dependency. The C port mirrors the Go firmware's data structures: a `firmware\_t` struct holds the complete state (routing table, MoE map, weight commands, KV cache, dispatch records) as statically sized arrays, with configuration constants that scale by compile target---under `PNM\_MCU` reduced limits for small microcontrollers (up to `PNM\_MAX\_LAYERS=2`, `PNM\_MAX\_NODES\_TOTAL=32`, `PNM\_MAX\_EXPERTS=16`, and `PNM\_KV\_CACHE\_DEPTH=64`), and the full SoC defaults otherwise (`PNM\_MAX\_LAYERS=8`, `PNM\_MAX\_NODES\_TOTAL=512`, `PNM\_MAX\_EXPERTS=256`, and `PNM\_KV\_CACHE\_DEPTH=4096`); it implements the identical five-phase boot and dispatch loop and compiles cleanly with `gcc -Wall -Wextra -std=c11` against the Go firmware's golden model.

The execution model is weight-stationary: every parameter and activation buffer resides in a fixed LPDDR6 bank at a known coordinate, and computation is triggered by data arrival (the doorbell of Section \hyperref[sec:2.9]{2.9}), not by instruction fetch. The MAC array reads weights from the local bank at the full 256 GB/s bus rate, executes the resident kernel, and returns the result through the fabric. There are no runtime page faults, no TLB misses, no cache coherence invalidations, and no operating-system context switches. The hardware is a physical dataflow graph; the software is the static assignment of that graph to silicon.

\subsection*{2.7 The compilation pipeline}\label{sec:2.7}

The compiler translates a model definition and its pre-trained weight tensors into a static dataflow graph that maps directly onto the physical chassis. The primary input is a safetensors archive---a flat binary of weight tensors with a JSON header indexing each tensor by name, dtype, shape, and byte offset---paired with a model configuration that specifies the computation graph: layer count, hidden dimensions, attention heads, expert counts, and activation functions. The output is a concrete assignment of weights, kernels, and routing tables to specific physical coordinates, plus a dispatch schedule that the central router chip executes at runtime. The backend does not target CUDA, SIMD, or a conventional instruction set. It targets the physical topography schema exported during POST (Section \hyperref[sec:2.5]{2.5}).

This is a fundamentally different compilation model than conventional GPU software. A CUDA compiler generates an instruction stream that a general-purpose core fetches, decodes, and executes against a shared memory hierarchy. The PNM compiler generates neither instructions nor a program counter: the output is a spatial assignment of weights and kernels to physical addresses, where the address IS the coordinate and the doorbell IS the program counter. Computation is triggered by data arrival, not by instruction fetch. The compiler's job is to make the static dataflow graph correct by construction, so that the hardware can execute it with zero runtime overhead.

The pipeline proceeds in five stages, each deterministic and side-effect-free.

Stage 1: Ingestion. The compiler reads two files from a HuggingFace model directory: the safetensors index (`model.safetensors.index.json`) and the model configuration (`config.json`). The index maps each named tensor to its shard filename; the config specifies the computation graph---layer count, hidden dimensions, attention heads, expert counts, and activation functions. Crucially, the safetensors index does not store tensor shapes: the compiler infers shapes from tensor names using naming conventions and the config parameters. For example, `model.layers.0.self attn.q proj.weight` maps to $[\mathrm{HiddenSize}, \mathrm{HiddenSize}]$; `model.layers.0.experts.3.gate\_up\_proj` maps to $[2 \times \mathrm{MoEIntermediateSize}, \mathrm{HiddenSize}]$; `model.layers.0.experts.3.down\_proj` maps to $[\mathrm{HiddenSize}, \mathrm{MoEIntermediateSize}]$; `model.embed\_tokens.weight` maps to $[\mathrm{VocabSize}, \mathrm{HiddenSize}]$; and `model.layers.0.router.proj.weight` maps to $[\mathrm{NumExperts}, \mathrm{HiddenSize}]$. Each tensor's byte length is computed as $\mathrm{params} \times 2$ for BF16 (2 bytes per element). The catalog records the name, shard, byte length, element count, and data type for every tensor. For a LLaMA-class MoE model, this yields on the order of 200--300 named tensors per layer (attention projections, feed-forward gates, layer norms, routing weights) plus shared tensors (embeddings, output head, global norms). The safetensors format is weight-only: it stores no computation graph, no gradient tape, and no optimizer state, which is exactly what inference compilation needs. The compiler reconstructs the graph from the config and the tensor naming conventions, not from the archive itself.

Stage 2: Capacity-aware partitioning. Each reference node holds 128 GB of LPDDR6 (Section \hyperref[sec:3.1]{3.1}). The partitioner assigns tensor shards to physical nodes subject to two constraints: (a) no node exceeds its 128 GB memory budget, and (b) tensors that participate in the same kernel invocation must be co-located or connected by a known route. The partitioning proceeds in three sub-phases.

Sub-phase 2a: Layer-to-physical-layer mapping. The compiler divides the model's hidden layers evenly across the chassis's physical layers: $\mathrm{model\_layers\_per\_physical} = \lceil \mathrm{NumHiddenLayers} / \mathrm{ChassisLayers} \rceil$. For a 30-layer model on a 4-layer chassis, each physical layer owns 8 model layers (the last owns 6). Physical layer $pl$ owns model layers $[pl \times 8, (pl+1) \times 8)$.

Sub-phase 2b: Expert distribution. Within each physical layer, the compiler distributes experts round-robin across the layer's $B_x \times B_y$ nodes. For model layer $ml$ and expert index $expIdx$, the target node is $nodeIdx = expIdx \bmod (B_x \times B_y)$, with coordinates $x = \lfloor nodeIdx / B_y \rfloor$, $y = nodeIdx \bmod B_y$. Each expert contributes two tensors: `gate\_up\_proj` of size $[2 \times \mathrm{MoEIntermediateSize}, \mathrm{HiddenSize}]$ and `down\_proj` of size $[\mathrm{HiddenSize}, \mathrm{MoEIntermediateSize}]$, both $\times 2$ bytes for BF16. A single expert's weights at BF16 for a 70B-class model occupy on the order of 70--80 GB, fitting one node with headroom for KV cache and activation buffers. Hot experts---identified from offline activation histograms---are replicated across $k$ nodes, placed on the same layer to avoid cross-layer spine traffic for the most frequently dispatched tokens (Section \hyperref[sec:2.12]{2.12}).

Sub-phase 2c: Dense components. Attention projections (Q, K, V, O), the dense MLP (gate, up, down projections), and two layer norms are assigned to a dedicated attention node, round-robin by model layer: $attnNodeIdx = ml \bmod (B_x \times B_y)$. This places all dense operations for a given model layer on a single node, eliminating cross-node traffic for the attention path. Embeddings are sharded across all layer-0 nodes by token-ID range: $\mathrm{rowsPerNode} = \lfloor \mathrm{VocabSize} / \mathrm{nodesLayer0} \rfloor$, with the remainder distributed to the first nodes. The final layer norm goes to a rotating layer-0 node: $normNodeIdx = \mathrm{NumHiddenLayers} \bmod \mathrm{nodesLayer0}$. Router gating weights ($router.proj.weight$) are assigned to a sentinel node at coordinate $(-1, -1, -1)$---the router chip itself, not a compute node.

Sub-phase 2d: KV cache reservation. On every node that holds attention Q tensors, the compiler reserves space for the KV cache: $\mathrm{kvEntryBytes} = \mathrm{HiddenSize} \times 4$ (K and V, each 2 bytes in BF16, per hidden dimension). The KV budget per node is 80\% of the 128 GB capacity, capped at 16,384 entries (a 16K context window). If the budget is exhausted, the firmware evicts the least-recently-used entries (Section \hyperref[sec:2.12]{2.12}).

For models whose individual tensors exceed node capacity---extreme tensor-parallel scenarios---the partitioner shards the tensor across multiple nodes on the same layer, and the compiler inserts a local reduction kernel that merges partial results via single-hop on-board routes that never enter the spine. This is a special case: most production MoE models fit comfortably within the per-node capacity at FP16 or INT8, which is the architectural sweet spot the memory sizing was designed for (Section \hyperref[sec:3.3]{3.3}).

Stage 3: Kernel mapping. Each tensor operation in the graph is mapped to a hardware kernel and a compute unit type. The mapping is role-based: MoE expert weight-streaming uses BF16 FMA (weight-stationary multiply), attention QKV projections use BF16 systolic MAC arrays, layer norm uses FP32 ALU (for the division in variance computation), and the gating network's softmax uses FP32 ALU MIN/MAX/CMP operations. The reference kernel set includes: dot (matrix-vector multiply, for linear projections), accum (running accumulation, for softmax denominators and layer norm variance), sum (element-wise addition, for bias and residual connections), and echo (pass-through, for validation nodes and pipeline staging). The DUV MAC ASIC's fixed-function array executes the kernel on data present in its local LPDDR6 bank. The compiler emits the kernel name, resident weights (as hex bytes in the .pnm directive format), and bias constant for each assigned node. For MoE routing weights, the central router chip holds the gating network and evaluates the closed-form coordinate function at runtime (Section \hyperref[sec:2.8]{2.8}); the per-expert kernels are static and pre-loaded.

The key insight is that the kernel set is deliberately small and fixed. There is no general-purpose instruction decoder, no branch predictor, and no speculative execution on the node. Each kernel is a pure function of its input payload and resident weights, which is what makes the doorbell mechanism (Section \hyperref[sec:2.9]{2.9}) sound: the hardware can assert DOORBELL\_TRIG with confidence that the node will produce a deterministic result, because there is no possibility of a page fault, a cache miss, or a branch misprediction changing the outcome. The compiler's kernel mapping is a many-to-one function---many tensor operations map to the same dot kernel, differing only in their weight dimensions---and the hardware executes them all identically.

Stage 4: Coordinate assignment and route computation. The partitioner's node assignments are translated into physical [Layer ID | Module ID] coordinates. The compiler then computes two static lookup tables.

The routing bitmap table: for each node at $(L, X, Y)$, the compiler computes an 11-bit bitmap: bits [10:7] hold the 1-based Layer ID ($(L+1) \ll 7$), bit [6] selects AXIS (0 = X), bit [5] selects SIGN (0 = +), and bits [4:0] hold the Y-distance from the lxy\_repeater. The lxy\_repeater at layer $L$ receives a bitmap matching Layer $L$; the node itself receives a bitmap matching its full coordinate with AXIS = Y and DIST = $Y$. These bitmaps are loaded into the fabric during POST (Section \hyperref[sec:2.5]{2.5}) and are immutable within a boot session.

The MoE expert map: for every tensor with role ``expert\_gate\_up,'' the compiler records a mapping from (model\_layer, expert\_idx) to the physical NodeID $(L, X, Y)$ that holds that expert's weights. At runtime, the router chip evaluates the gating network, selects the top-$k$ experts, and looks up this table to determine the destination coordinates. The table is static and loaded into the router's SRAM during Phase 4 of POST.

Because the fabric is a single-spine tree with dimension-order on-board routing, there is exactly one legal route between any pair of coordinates, and the route is an $O(1)$ arithmetic function of the two coordinate pairs (Section \hyperref[sec:2.8]{2.8}). No pathfinding algorithm is invoked. The compiler also computes the deterministic latency for every route: hop count times per-hop delay plus message serialization at the link rate (Section \hyperref[sec:3.4]{3.4}). These latencies become the compile-time budgets that sources use to detect failures (Section \hyperref[sec:4.1]{4.1}). The complete routing table is a static, immutable data structure loaded into the chassis during POST (Section \hyperref[sec:2.5]{2.5}).

Stage 5: Dispatch schedule generation. The AOT mapper produces a static execution schedule that assigns injection timestamps to every message in the workload. The schedule respects three invariants: (a) no node's elastic input buffer (Section \hyperref[sec:4.3]{4.3}) is overcommitted, (b) every spine link's aggregate bandwidth stays within the sizing rule of Section \hyperref[sec:2.1]{2.1}, and (c) every message completes within its closed-form latency budget. For MoE workloads, the schedule is parametric in the token-to-expert mapping---the router chip fills in the concrete coordinates at runtime via the expert map, but the schedule guarantees that any legal mapping satisfies the buffer credits. For stencil workloads, the schedule is fully static: halo exchanges, compute phases, and result writes are all pre-timed.

The concrete compilation path from a safetensors archive to executing hardware follows a straight pipeline. The ingest stage produces a tensor catalog and computation graph from the safetensors header and model config. The partition stage assigns weight shards to node coordinates via layer mapping, expert round-robin, dense node assignment, and embedding sharding. The map stage assigns operations to kernel names and compute unit types based on tensor roles. The route stage computes routing bitmaps and the MoE expert map. The schedule stage produces a static dispatch timeline. The emit stage writes the .pnm program file and weight binary images. The load stage streams weights to nodes during POST (Section \hyperref[sec:2.5]{2.5}). The execute stage has the router dispatch tokens and the doorbells fire.

The .pnm program format is the concrete realization of this pipeline: each directive (kernel, bias, token) corresponds to a specific stage's output. A kernel directive encodes the map stage's kernel assignment and the partition stage's weight shard. A bias directive encodes the MAC stub's KERNEL\_CONST constant, which the hardware applies to every arriving payload before the resident kernel executes. A token directive encodes the schedule stage's injection: the destination coordinate, the payload bytes, and the implicit routing that the route stage computed. The co-simulation harness (Section \hyperref[sec:2.15]{2.15}) exercises this path end-to-end: the pnmc compiler lowers a plain-text program onto the chassis, generates Verilog stimulus, compiles the fabric slices, simulates them, and verifies hardware-delivered results against the golden model.

A concrete example of the .pnm format. A two-expert program targeting a 512-node chassis: the first expert is a byte-sum kernel on node (board 3, column 2, row 5) with a silicon bias of +7, and the second is a dot-product kernel on node (board 8, column 8, row 8) with bias +3 and a 16-byte resident weight vector. Five tokens are dispatched to both nodes. The .pnm directives encode this as:

kernel sum 2 2 5
bias   7 2 2 5
kernel dot 7 7 7 01 02 03 04 05 06 07 08 09 0a 0b 0c 0d 0e 0f 10
bias   3 7 7 7
token 2 2 5   10 20 30 40 50 60 70 80 90 a0 b0 c0 d0 e0 f0 00
token 2 2 5   01 02 03 04 05 06 07 08
token 7 7 7   de ad be ef
token 0 0 0   11 22 33 44 55 66
token 4 3 1   a5 5a 99 00 ff

Each directive line specifies the kernel name, destination coordinates, and---for kernel directives---the resident weight bytes as hex. The compiler lowers these directives into the wire format of Section \hyperref[sec:2.9]{2.9}: each token becomes a flit carrying [LAYER\_ID | MODULE\_ID | CTRL | LEN\_LO | LEN\_HI | payload | CRC\_HI | CRC\_LO], and the co-simulation harness verifies byte-exact delivery and kernel results against the golden model.

A working implementation of this pipeline exists in the companion repository: the `pnmc` compiler accepts a HuggingFace model directory (containing `config.json` and `model.safetensors.index.json`), parses the safetensors header to extract tensor metadata, computes shapes from the model configuration, and performs the five-stage compilation. For example, compiling Google's Gemma-4-26B-A4B-it (a 30-layer MoE model with 128 experts, 8 active per token, 26.5B parameters in BF16) onto a 4-layer, 4$\times$4 = 64-node chassis assigns 8 model layers per physical layer, distributes 128 experts round-robin across 16 nodes per layer, and produces a chassis schema with coordinate maps, routing bitmaps, and a .pnm program compatible with the co-simulation harness. The compiler is model-agnostic: it reads any safetensors index and config combination and maps it onto any chassis dimensions.

The compilation model has a second, less obvious consequence: because every weight is statically assigned to a fixed coordinate, and every route is a deterministic function of that coordinate, the compiler can reason about the entire machine's state at compile time. There are no runtime page faults, no TLB misses, no cache coherence invalidations, and no operating-system context switches. The weight loading phase (Section \hyperref[sec:2.5]{2.5}) is a one-time DMA burst that fills each node's LPDDR6 bank; once loaded, weights never move. When a token arrives at a node, the doorbell fires, the MAC reads weights from the local bank at the full 256 GB/s bus rate, and the result is returned. The compiler has already verified that the weight dimensions fit the node, that the route is valid, that the latency budget is sufficient, and that no buffer will overflow. The hardware merely executes what the compiler proved correct.

The source-language front ends (Section \hyperref[sec:2.6]{2.6}) feed into this pipeline: the R and Haskell compilers emit FP64 IR that the kernel mapper assigns to fp64\_fma units, while the HLSL compiler emits FP32 ALU IR that the mapper assigns to fp32\_alu units. The mapper also handles BF16 and FP16 FMA operations for transformer inference (Section \hyperref[sec:2.12]{2.12}), INT8 MAC operations for quantized inference, and systolic array operations for attention QKV projections. The kernel set is deliberately small and fixed: each kernel is a pure function of its input payload and resident weights, which is what makes the doorbell mechanism (Section \hyperref[sec:2.9]{2.9}) sound. The compiler's kernel mapping is a many-to-one function---many tensor operations map to the same kernel type, differing only in their weight dimensions---and the hardware executes them all identically.

\subsection*{2.8 Coordinate routing and dispatch}\label{sec:2.8}

The compilation pipeline operates in two phases.

Ahead-of-Time (AOT) mapping. The compiler assigns weights, kernels, and grid subdomains to physical coordinates, then computes every route and latency in closed form: within a layer, X-then-Y between coordinates; across layers, a monotone spine traversal between layer attachments, bracketed by on-board paths. Both are $O(1)$ arithmetic functions of the [Layer ID | Module ID] pair---on a single-path tree there is exactly one route, so no pathfinding algorithm has anything to find. The host-side driver orchestrates this mapping: it pre-computes 11-bit routing bitmaps for every node (bits [10:7] hold the 1-based Layer ID shifted left by 7, bits [4:0] hold the Y-distance from the lxy\_repeater), builds a MoE expert map that associates each (model\_layer, expert\_idx) pair with the physical NodeID holding that expert's weights, and generates the weight upload commands sorted lexicographically by $(L, X, Y)$ for deterministic ordering. Each upload command carries the target Layer ID, Module ID (the byte $X[3:0] \parallel Y[3:0]$), payload length, weight bytes, the compute unit type, and the data type, enabling the host to verify weight upload integrity per unit type. The driver also computes the 11-bit routing bitmap for the router chip itself, which is assigned a sentinel coordinate $(-1, -1, -1)$ and is the only node not loaded into the fabric's routing tables.

Just-in-Time (JIT) dispatch. Dynamic tokens---such as MoE expert routing decisions---arrive at the central router chip. The router evaluates the same closed-form routing function, prepends a strict [Layer ID | Module ID] header, and injects. There is no runtime arbitration, no congestion sensing, and no adaptive routing. The router's 1 GHz dispatch stage issues one routing decision per clock, but the sustained dispatch rate is set by the fabric, not the decision stage: injection is bounded by the destination buffer credit rule of Section \hyperref[sec:4.3]{4.3}, cross-layer traffic by the spine sizing rule of Section \hyperref[sec:2.1]{2.1}, and per-link traffic by the 16 GB/s link rate. At the reference spine of Section \hyperref[sec:2.1]{2.1} the sustained cross-layer dispatch rate is therefore on the order of $10^{8}$--$10^{9}$ dispatches/s depending on expert placement---corresponding, at 270 dispatches per generated token (Section \hyperref[sec:2.6]{2.6}), to a token rate of roughly $10^{6}$ tokens/s---comparable to what 131 TB/s of node-local memory can serve, and far below the router's raw decision capacity. Optional per-layer dispatch replicas scale the decision rate linearly and remain deterministic because the routing function is pure; they do not raise the fabric-bound ceiling.

\subsection*{2.9 The doorbell mechanism}\label{sec:2.9}

A flit's lifecycle exemplifies zero-overhead activation. The central router chip injects the payload into the fabric. The combinational comparator inside the lxy\_repeater at the target layer gates the flit onto the X/Y NoB in sub-nanosecond time. Hardware Flit Repeaters pass the byte stream via DMA into the target node's local address space. The message header carries a length field and a CRC-protected destination coordinate; the node's DMA engine---the memory controller's DMA engine, Section \hyperref[sec:2.16]{2.16}---counts incoming bytes, computes the running message CRC over the protected body, and compares the protected destination to its own coordinate. When the byte count equals the length field, the CRC validates, and the destination matches, the DMA engine asserts DOORBELL\_TRIG and DOORBELL\_ACK. This physical pin transition wakes the DUV MAC ASIC instantly; there is no software interrupt handler, no software polling loop, and no operating-system context switch.

The hardware fabric decouples transport from computation. The deterministic repeater network connects to a modular Processing Element (PE) interface via standard AXI-Stream / Ready-Valid handshaking, allowing the underlying execution block (systolic array, FP16 MAC array, or ALU) to be swapped independently of the routing fabric. The companion PE tile stub in the HDL sketch (HDL/core/pe\_tile\_stub.v) implements this contract: it consumes a flit on its slave port, holds it for one to two clock cycles to model multiply-accumulate latency, and re-presents it on the master port, so the transport-level verification of Sections \hyperref[sec:3.4]{3.4} and \hyperref[sec:4.3]{4.3} is independent of which execution block sits behind the interface. The stub also computes its resident kernel in silicon---an element-wise bias add---and recomputes the end-to-end CRC over the transformed body, so the co-simulation (Section \hyperref[sec:2.15]{2.15}) cross-checks the hardware verdict against the software twin on every message.

Every node---whether a MAC node or a repeater serving a CAMM slot---runs the same hardwired three-phase doorbell discipline. This is not software polling: no instructions are fetched and no program counter exists. The phases are the armed and drain states of a finite state machine burned into the doorbell logic:

WATCH (initial phase): continuously check the sentinel bit at the end address of the DMA target buffer. When the sentinel is set (byte count equals header length, CRC valid, and protected destination matches), interrupt WATCH and enter DISPATCH.

DISPATCH: fire the node's resident function. COMPUTE (PNM processor): execute the resident kernel, e.g., matrix multiplication over the landed payload and local weights. FORWARD (repeater at the CAMM slot): stream the payload onward to the next statically mapped node. DOORBELL\_ACK is asserted with DOORBELL\_TRIG at the fire edge.

DRAIN (reverse phase): while results egress, check the sentinel; when transfer completes, clear the sentinel (bit now empty) and clear the initial payload data; restart WATCH.

Result egress is the same ASIC, not a separate chip: after the kernel completes, the transmit DMA---the reverse half of the doorbell's DMA engine---reads the result from the MAC staging buffer, assembles a return flit whose DEST is the requester (a coordinate the AOT schedule fixes at compile time, Section \hyperref[sec:2.8]{2.8}), appends the end-to-end CRC, and streams it out the AXI-Stream master port into the fabric (Table 4); the fabric's stateless gates then forward it exactly like any other flit, so the return path adds no new silicon.

The buffer is rigorously empty before the next message can arrive, making back-to-back activations safe without software synchronization. Because every route is a single path of strictly in-order HFR stages, byte-count equality is a sound completion test. An aborted stream always fails the count, the CRC, or the destination check, so the doorbell never fires on a partial, corrupt, or misdelivered payload---and when it refuses, DOORBELL\_ACK is withheld, which the source's watchdog converts into an observable NODE\_ERR (Section \hyperref[sec:4.1]{4.1}).

\subsection*{2.10 Reliability and determinism}\label{sec:2.10}

LPDDR6 provides on-die ECC; the MAC ASIC's memory controller adds conventional SECDED on the bus. Every fabric link carries a link-local CRC per flit that covers the full flit including its routing header bytes, so a bit error that would corrupt a route decision is detected before the consuming gate strips its header byte; a corrupt flit is discarded at the boundary rather than misrouted. Every message additionally carries an end-to-end CRC per message over the protected body---[CTRL, DEST, LEN, payload]---and the destination DMA refuses the doorbell unless the protected destination matches its own coordinate (Section \hyperref[sec:2.9]{2.9}).

Deterministic replay is claimed precisely, in two distinct senses. First, re-injecting an identical stimulus onto an identical fabric reproduces delivery bit-exactly: the fabric has no hidden state, and this property is machine-checked in the co-simulation (Sections \hyperref[sec:3.5]{3.5} and \hyperref[sec:4.3]{4.3}). Second, DRAM contents are not assumed to be bit-stable across arbitrarily long runs: LPDDR6 on-die ECC plus the bus SECDED protect every access, and a background scrub pass re-reads resident weights at a rate sized so the mean time to an undetected mis-correction exceeds the chassis lifetime. Post-fault recovery (Section \hyperref[sec:4.1]{4.1}) re-runs the work on a healthy replica; that is functional recomputation under the same doorbell discipline, not bit-exact replay of the failed node's memory contents. The machine guarantees determinism of the fabric and of computation given a memory snapshot; soft errors in memory are handled statistically by scrub and retry, not claimed away.

\subsection*{2.11 Target Workloads}\label{sec:2.11}

\subsubsection*{2.11.1 Workload simulation suite}\label{sec:2.11.1}

Five canonical workloads exercise distinct routing patterns against the co-simulation harness of Section \hyperref[sec:2.15]{2.15}, serving as both correctness tests and routing-pattern demonstrations. Each is generated as a .pnm program (Section \hyperref[sec:2.7]{2.7}) via the pnmc workload command and verified byte-exact on the gate-level fabric.

jacobi5 maps a 5-point Jacobi stencil onto a single physical layer's X/Y grid: each node holds one grid point with stencil coefficients [w\_c, w\_n, w\_s, w\_e, w\_w] as resident dot-kernel weights in Q6 fixed-point, and each token packs the 5-byte neighborhood [center, north, south, east, west]. Routing is pure intra-layer dimension-order X$\rightarrow$Y---zero spine traffic, exactly matching the halo-exchange pattern of Section \hyperref[sec:2.13]{2.13}.

matvec distributes a matrix row-per-node across all layers: each node's resident dot kernel carries one row's weights (deterministic entries $(i+j) \bmod 251 + 1$), and the input vector is dispatched to every row-node simultaneously. Routing uses spine descent for cross-layer rows; within each layer, X$\rightarrow$Y dimension-order. This is the weight-stationary pattern that Section \hyperref[sec:3.2]{3.2}'s roofline analysis targets.

reduction exercises reverse-path merge arbitration: leaf nodes on layer 0 send data fragments toward a root node at the top layer, traversing board egress $\rightarrow$ Y-up $\rightarrow$ X-up $\rightarrow$ lxy\_repeater $\rightarrow$ spine ascent. All nodes run the stateless echo kernel so the golden model is order-independent regardless of merge-point round-robin decisions (a production reduction tree would use accum at the root); what the workload proves is that the arbitrated egress path delivers every packet intact under contention.

broadcast distributes identical payloads from the root to all nodes---the POST Phase-3 weight-upload pattern (Section \hyperref[sec:2.5]{2.5}) at reduced scale. Routing follows spine descent to every layer, then per-layer X$\rightarrow$Y fan-out; it stresses spine-descent bandwidth and the worst-case hotspot on the source node's uplink.

nbody generates all-pairs O($N^2$) traffic: every node sends a token to every other node. This saturates every link class simultaneously and serves as an upper-bound benchmark only, capped at 64 nodes for simulation tractability.

All five pass the full verification suite (byte-exact delivery, doorbell accounting, kernel correctness, latency bounds) on the default chassis.

\subsubsection*{2.11.2 Routing designs by use case}\label{sec:2.11.2}

The fabric supports four routing patterns, selected by the AOT compiler based on the algorithm's communication topology rather than programmed at runtime:

Dimension-order X$\rightarrow$Y (intra-layer). Nearest-neighbor stencils, halo exchanges, and any communication confined to a single board. The route never enters the spine, so its latency is a compile-time constant independent of chassis depth: a single dimension-order hop for a nearest-neighbor halo exchange, plus PE pipe delay. This is the cheapest pattern and should be preferred whenever data locality permits.

Spine descent (cross-layer unicast). Weight-stationary linear algebra (matvec), MoE expert dispatch, and any pattern where the injector reaches nodes on deeper layers. Each hop adds exactly one spine stage; the closed-form latency of Section \hyperref[sec:3.4]{3.4} accounts for this deterministically.

Reverse-path merge (egress collection). Reduction trees, result aggregation, and any many-to-one pattern. The forward path is a demux tree making no arbitration decisions; the reverse path is a pure merge tree whose round-robin grants are held for whole packets (never torn) and whose dependency graph is acyclic, so deadlock-free delivery is guaranteed structurally (Section \hyperref[sec:4.3]{4.3}).

Spine descent + fan-out (broadcast/multicast). Weight upload, parameter distribution, and any one-to-all or one-to-many pattern. The source node's uplink is the bottleneck; the compiler schedules broadcast bursts during idle phases (POST boot) or overlaps them with compute to hide serialization.

The compiler selects among these patterns during its route-planning stage: it computes each message's source-destination pair, checks whether the pair shares a layer (dimension-order suffices), and otherwise inserts spine hops into the source-routed header. No runtime routing decisions exist in the fabric itself---every flit carries its complete route from the moment of injection.

\subsection*{2.12 Mixture-of-Experts transformers}\label{sec:2.12}

Trillion-parameter MoE transformers present a capacity crisis for HBM-based systems: only a sparse subset of experts is active per token, yet the entire parameter set must be accessible\cite{3,4}. Our architecture distributes expert weights across hundreds of cheap LPDDR6 pools. A one-trillion-parameter model at INT8 occupies 1 TB---eight reference nodes---and at FP16, sixteen. A single 64 TB chassis holds tens of trillions of parameters with room for replication.

When a token requires a given expert, the central router chip evaluates the coordinate routing function, dispatches the token via wormhole routing to the holding node, and the doorbell fires. The local MAC executes and returns the updated hidden state. Because only the token travels to the expert---the expert's weights stay resident---the cross-layer traffic fraction $\alpha$ of Section \hyperref[sec:2.1]{2.1} is a direct consequence of expert placement: with random placement across eight layers, $\alpha = 7/8$ of token dispatches cross the spine, and the honest spine-sizing rule sizes the fabric for exactly that. Three refinements matter in production: (i) hot-expert replication from offline activation histograms across $k$ nodes, which both cuts latency and---when replicas are placed layer-locally---shrinks $\alpha$ toward the resident dense-component fraction, trading replication cost against spine provisioning; (ii) expert capacity via elastic input buffers sized to at least one full message burst, with compiler-set capacity factors; (iii) dense components (attention projections, embeddings, output head) replicated on dedicated dense nodes per layer under the same doorbell discipline. Idle experts draw no dynamic compute power---only DRAM refresh current.

The compute unit assignment for MoE inference follows a fixed mapping: expert weight-streaming uses BF16 FMA (weight-stationary multiply), attention QKV projections use BF16 systolic MAC arrays, layer normalization uses FP32 ALU (for the division in variance computation), and the gating network's softmax uses FP32 ALU MIN/MAX/CMP operations. The model compiler (Section \hyperref[sec:2.7]{2.7}) performs this assignment during the kernel mapping stage, and the firmware (Section \hyperref[sec:2.6]{2.6}) records the unit type in each dispatch record so the host-side driver can verify weight upload integrity per unit type. The reference chassis sustains 30 dense dispatches plus 240 MoE dispatches per generated token at the dispatch rates of Section \hyperref[sec:2.8]{2.8}, with the BF16 FMA units consuming approximately 2 bytes per FLOP at the 256 GB/s memory bandwidth.

Expert selection itself is not end-to-end trained in this work. The co-simulated dispatch path routes every token through a deterministic token-seeded routing hash---a 64-bit FNV-1a seed over the token bytes and model layer, mixed per expert, reproducible run-to-run for a fixed input---while the RTL's moe\_gating unit (Section \hyperref[sec:2.17]{2.17}) provides the learned-projection datapath (hidden $\times$ router\_weights, top-$k$ selection) verified structurally by testbench with hand-loaded weights. The two paths share the selection shape but not the learned behavior: the hash is a demonstrator for dispatch mechanics, and a learned $128$-expert BF16 projector at the gating dimension of the reference model would need roughly $0.3$--1 MB---5--15$\times$ the router's 64 KB gating SRAM---so the learned-router sizing of Section \hyperref[sec:2.5]{2.5} and its per-token evaluation latency are projections (Section \hyperref[sec:4.4]{4.4}).

The KV cache is distributed across the attention nodes: on every node that holds attention Q tensors, the compiler reserves 80\% of the 128 GB capacity for KV entries, capped at 16,384 entries (a 16K context window). Each KV entry occupies $\mathrm{HiddenSize} \times 4$ bytes (K and V, each 2 bytes in BF16, per hidden dimension). Entries are distributed round-robin across four directional banks (X+, X-, Y+, Y-) by sequence position, and eviction is triggered when any bank exceeds the 80\% capacity threshold. The firmware checks the KV cache after each attention dispatch: if a bank needs offloading, it dispatches a KV offload command through the spine to host DRAM and triggers a reload cycle before the next token. The KV cache is entirely hardware-managed---no software polling, no OS page faults---and its FIFO eviction policy is deterministic given the same sequence of token arrivals.

For interactive inference (chat, code generation, retrieval-augmented generation), the architecture includes an LLM inference client that handles tokenization, prefill, autoregressive generation, and temperature/nucleus sampling. The client boots the firmware, builds a vocabulary from the model's tokenizer configuration, and generates tokens one at a time: each token triggers a prefill phase (all layers, all experts) followed by an autoregressive phase (cached KV state, reduced compute). The sampling pipeline computes softmax over the logit vector with numerical stability (max-subtraction trick), applies temperature scaling, sorts by probability, and selects via nucleus (top-p) sampling, all running on the host with the results dispatched to the chassis as token directives. The client tracks per-inference statistics---tokens generated, prefill tokens, total dispatches, MoE and dense dispatches, KV store/load/eviction operations---and reports compute unit utilization across the chassis. Token encoding on the wire varies by data type: 2 bytes for BF16 and FP16, 4 bytes for FP32, and 8 bytes for FP64.

\subsection*{2.13 Scientific HPC stencil workflows}\label{sec:2.13}

Climate modeling, computational fluid dynamics, and seismic imaging rely on regular FP64 stencil operations over massive grids. These workloads map onto the rigid X/Y grid of sockets: each grid subdomain becomes a physical node, and halo exchanges become single-hop dimension-order routes between neighboring sockets that never enter the spine. Because the stencil neighborhood is known at compile time, the AOT compiler generates exact routes with no runtime indirection, replacing cluster middleware such as MPI\cite{19} with a coordinate function.

The roofline analysis of Section \hyperref[sec:3.2]{3.2} shows why mature nodes suffice: a 7-point FP64 stencil has an arithmetic intensity of roughly 0.3--0.5 FLOP/byte, far below the reference node's machine balance of about 1.5 FLOP/byte\cite{5}. Where numerics permit, mixed-precision halo storage raises effective intensity further. Deterministic worst-case halo-exchange latency is a compile-time constant (Section \hyperref[sec:4.4]{4.4}).

\subsection*{2.14 Functional data-based workflows}\label{sec:2.14}

The mappable subset is first-order dataflow graphs of pure kernels over statically allocated buffers\cite{6,11}. Within this subset, an actor is a node, a message is a flit, and a statically scheduled functional kernel is a doorbell activation. General lazy graph reduction with an unbounded dynamic heap is future work (the Conclusion section). Referential transparency is exactly what makes the replay guarantees of Sections \hyperref[sec:2.6]{2.6} and \hyperref[sec:4.1]{4.1} hold.

\subsection*{2.15 Verification methodology: machine-checking the transport claims}\label{sec:2.15}

The transport claims of this paper are machine-checked, not asserted. The companion repository ships a co-simulation harness---Go, standard library only---that drives a cycle-exact Verilog-2005 model of the fabric: the routing gates of Sections \hyperref[sec:2.1]{2.1} and \hyperref[sec:2.2]{2.2}, the per-node PE tile stub of Section \hyperref[sec:2.9]{2.9}, and the end-to-end CRC machinery of Section \hyperref[sec:2.10]{2.10}\cite{28}. The harness is deterministic by construction: it embeds a byte-exact MT19937 clone of the CPython random.Random generator (a Python-equivalent Mersenne Twister implementation verified against CPython's known outputs), so a given seed reproduces the identical stimulus file, generated Verilog, and delivery log on any machine, and every result quoted in Section \hyperref[sec:3.5]{3.5} is bit-reproducible. A software oracle checks five properties on every run: (1) byte-exact delivery and byte conservation---the bytes delivered to each node's DMA port equal the bytes injected for it, and the totals match, so no byte is lost, duplicated, or reordered anywhere in the fabric; (2) zero drops and zero misroutes---every injected flit reaches exactly its declared destination, and the routing-bitmap comparators of Section \hyperref[sec:2.1]{2.1} never raise route\_err on valid traffic, even under backpressure; (3) doorbell accounting---the hardware verdicts (activations, refusals, and the stub's corrupt\_out pulses) match, message by message, the software twin's three-condition fire logic of Section \hyperref[sec:2.9]{2.9}; (4) kernel correctness---each resident kernel (echo, sum, accum, dot) executes on what the hardware actually delivered and must equal the golden model computed over the payload, including the bias the MAC stub already added in silicon; (5) latency bounds---the sweep scenario checks every packet against the closed form of Section \hyperref[sec:3.4]{3.4} exactly (hop count times per-hop delay plus serialization at the link rate, plus the stub's elastic pipe), while the other scenarios assert the pipe floor only, because contention is a compile-time-scheduled term.

Six scenarios exercise the fabric. sweep sends one flit to every node and checks the exact closed-form latency; vcsweep repeats sweep while also injecting pass-through traffic on all four virtual-channel classes of Section \hyperref[sec:4.3]{4.3}, so a single run exercises the entire VC stack including the arbitrated egress tree; load sends 500 random flits (450 routed to node kernels plus 50 pass-through) to random kernels with 100/50/25\% sink backpressure and proves lossless delivery with slow consumers; hotspot concentrates MoE hot-expert traffic (Section \hyperref[sec:2.12]{2.12}) on a worst-case corner node behind a 1-in-8 DMA gate; stress injects roughly 24 flits per node with 4 B--1 KB payloads, corrupts the end-to-end CRC of approximately 3\% of messages---which the doorbell must refuse---and passes approximately 5\% through the fabric; replay runs stress twice and requires the two delivery logs to be bit-identical, machine-checking the no-hidden-state claim of Section \hyperref[sec:2.10]{2.10}. Three hand-written HDL testbenches complement the harness: a functional fabric smoke test, a 500-packet load test with backpressure and per-packet checksums, and a doorbell test that drives eight packets through the PE tile stub and checks the hardened fire conditions---six activations, two refusals (a truncated message and a wrong destination), and two stub corrupt\_out pulses on messages whose CRC fails in flight. Additional testbenches verify every compute unit type in isolation: FP16, BF16, FP32, and FP64 FMA modules (3-cycle pipeline, bit-exact golden-model checks), the FP32 ALU (restoring binary long division, MIN/MAX/CMP, and FMA-path operations), the INT8 MAC (two-cycle multiply-accumulate), the BF16 and FP16 systolic MAC arrays (weight loading, burst activation, zero-weight identity), the MoE gating unit (softmax top-k selection), and the router chip (POST discovery, routing-table load, weight upload, MoE map load, and inference dispatch). Beyond the fabric and compute silicon, tb\_sodimm\_ctrl.v verifies the banked memory controller described in Section \hyperref[sec:2.17]{2.17}---exact row-hit/miss/conflict latencies, tRAS enforcement, byte-enable masking, refresh-storm integrity, and a DDR4-versus-DDR5 ordering sweep---and tb\_optical\_link.v exercises the optical inter-chassis links of Section \hyperref[sec:2.19]{2.19}: idle-link latency, burst framing alignment, flood-level losslessness, gapped packets, an asynchronous 100--80 MHz clock crossing, and simultaneous bidirectional traffic. A Verilator static lint pass confirms the routing fabric is free of multi-top errors and width mismatches across all gate-level modules. The AOT compiler path is exercised end-to-end by a small program compiler that lowers a plain-text program (kernel, bias, and token directives) onto the chassis and runs the identical pipeline; the shipped example places a byte-sum and a dot-product expert with silicon biases on the 512-node chassis and verifies their results on the gates (Section \hyperref[sec:3.5]{3.5}).

The forward path is a demux tree and makes no arbitration decisions, but the reverse path---node TX to Y-up, X-up, lxy\_repeater, and up-spine---is a pure merge tree whose arbitration is load-bearing (Section \hyperref[sec:4.3]{4.3}). Every merge is a 2-in/1-out module that grants one upstream port for a whole packet---the decision is held until EOP, so packets are never torn---and round-robins between the two ports; the granted port's ready tracks the downstream handshake while the other stalls in place. Because the merge dependency graph is a DAG (board egress toward spine ascent, never a cycle), arbitration cannot deadlock. The harness therefore validates the egress path twice: in the cycle-exact RTL itself, and through an independent activity-driven discrete-event model (DES) of the merge tree that shares no code with the Verilog. A cross-check oracle compares the DES delivery stream byte-for-byte against a real RTL slice on a mixed four-class program (Section \hyperref[sec:3.5]{3.5}), so agreement between the two implementations pins the arbitration logic down rather than echoing one implementation against itself.

vvp is single-threaded, so the harness partitions the layer list into contiguous groups---one per CPU core by default---and runs one vvp process per slice in parallel, compiling each slice's Verilog concurrently. The split is byte-exact against a monolithic simulation because the spine stages upstream of a slice's destination layers are transparent pass-through; cross-group spine contention is a compile-time scheduling concern, not a gate-level effect, and is deliberately left uncoupled (the AOT schedule of Section \hyperref[sec:2.8]{2.8} assigns it statically). Each slice's testbench injects from a stimulus file, applies per-node backpressure (each DMA sink accepts one byte in N cycles, N fixed per scenario), and terminates adaptively: it finishes 200 cycles after the last transition on any fabric wire, with a watchdog at 64 times the injected byte count for pathological stalls.

\subsection*{2.16 The LPDDR6 memory controller}\label{sec:2.16}

The memory controller is the ASIC's only memory-facing state machine, and it is the block that turns the physical bus of Section \hyperref[sec:2.4]{2.4} into the $\approx$256 GB/s figure quoted throughout this paper. It sits between the LPDDR6 PHY and the two consumers on the node---the MAC array and the doorbell DMA engine---and it is hardwired, not programmable: its policy is a fixed function of its inputs, which is what lets the compiler's static schedule (Section \hyperref[sec:2.8]{2.8}) treat memory service as a deterministic resource. Its place in the ASIC, pin group, and power budget are tabulated in Section \hyperref[sec:3.6]{3.6} (Table 4).

Bus structure. The 192-bit interface is four 48-bit channel groups, i.e. eight 24-bit LPDDR6 channels of two 12-bit sub-channels each---sixteen independent sub-channels on one module. Each sub-channel owns its own bank array, command path, and strobes, so the controller can keep up to sixteen rows open simultaneously and interleave transfers across them. This independence is the entire trick: with sixteen sub-channels in flight, the controller hides the two inherently serial DRAM operations---row activation and precharge---behind transfers on the other fifteen, sustaining full bus rate while no individual bank is ever busy more than a fraction of the time.

Access flow. The dominant operation is a long read burst: the MAC array consumes resident weights at the bus rate, and the AOT compiler lays weights out contiguously (Sections \hyperref[sec:2.7]{2.7} and \hyperref[sec:2.8]{2.8}), so each row activation serves a long run of column reads and the row-buffer hit rate approaches 100\% in steady state. The controller's bank scheduler issues activate, read/write, and precharge commands to the sub-channels in rotation, batching the residual random accesses---KV-cache updates and result writes---into write slots between weight streams. A bounded reordering window is safe because the fabric's closed-form latency bounds terminate at the node's DMA port (Sections \hyperref[sec:3.4]{3.4} and \hyperref[sec:4.3]{4.3}); kernel completion timing is not part of the verified bound, and results are byte-exact regardless of scheduling. The elastic input buffer of Table 4 (one or more maximum-length messages) decouples the MAC's demand from DRAM timing on both sides of the pipeline.

Refresh and ECC. The controller reserves statically sized refresh slots in the schedule---deterministic by construction, so the compiler can size the elastic buffers to cover them---and the isothermal manifold of Section \hyperref[sec:2.3]{2.3} keeps the DRAM at $\leq$70$^\circ$C, below the high-temperature refresh regime. Error protection is layered exactly as Section \hyperref[sec:2.10]{2.10} specifies: LPDDR6 on-die ECC corrects single-bit errors inside the DRAM, and the controller's SECDED code over the bus catches everything between the die and the ASIC, computed inline in the PHY datapath with no added pipeline latency. The background scrub pass of Section \hyperref[sec:2.10]{2.10} issues periodic read-and-verify cycles through the same bank scheduler at low priority; a scrub read never preempts a kernel's weight stream, and its rate is sized so the mean time to an undetected mis-correction exceeds the chassis lifetime.

Doorbell DMA path. The doorbell's DMA engine of Section \hyperref[sec:2.9]{2.9} lives inside the controller. Incoming flits are written into the target buffer as they arrive through the write path; the DMA engine maintains the byte count, computes the running CRC over the protected body, and compares the protected destination---all inline in the controller datapath. When the count reaches the length field, the CRC validates, and the destination matches, the controller sets the sentinel bit at the end address in the same cycle the doorbell fires (Section \hyperref[sec:2.9]{2.9}). The WATCH-phase sentinel check is a single read issued through the normal scheduler; it costs one bank access every few hundred cycles, negligible against the weight stream. The same write path loads weights at boot: the POST sequence of Section \hyperref[sec:2.5]{2.5} streams 64 TB through the router uplinks in under ten minutes as a long, purely sequential write burst per node---the controller's easiest workload.

Determinism. The controller holds timing state (refresh counters, bank timers, scheduler pointers) but no data state, and its policy is deterministic given its inputs, so it introduces no non-determinism into delivery or results: identical stimuli onto an identical chassis reproduce identical computation given a memory snapshot, exactly as Section \hyperref[sec:2.10]{2.10} claims. Residual DRAM timing variation is absorbed by the elastic buffers and never perturbs the data path, and soft errors are handled statistically by ECC, scrub, and the recompute protocol of Section \hyperref[sec:4.1]{4.1}---not claimed away.

Generation portability. The controller's bus contract---valid/ready acceptance, CAS-latency wait, one response per transaction---is deliberately the only interface its consumers see, which makes it portable across DRAM generations. Three behavioral PHY variants in the repository pin this down: an LPDDR5 profile (four banks, CAS 16, tRCD 4, tRP 3, tRAS 8, 3.9 $\mu$s refresh), an LPDDR5X profile (CAS 14, tRCD 3, tRP 2, tRAS 7), and the flagship LPDDR6 flat-bank profile (CAS 4). Each is exercised by a self-checking testbench that polls the ready strobe rather than assuming fixed delays (tb\_lpddr5\_phy.v, tb\_lpddr5x\_phy.v, tb\_lpddr6\_phy.v). A deployment selects the module class by synthesis-time parameter; the MAC array, doorbell DMA, and firmware are untouched---the same generation-agnostic claim Section \hyperref[sec:3.7]{3.7} machine-checks for SODIMM classes, extended to the LPDDR5/5X/6 CAMM2 ladder.

\subsection*{2.17 The BMC/Router chip: a RISC-V SoC for chassis management}\label{sec:2.17}

The central router chip described in Section \hyperref[sec:2.5]{2.5} is implemented as a RISC-V System-on-Chip (SoC) targeting Linux and Redox OS compatibility on bare-metal or NOMMU configurations. The SoC integrates an RV32IMA multi-cycle CPU core, 64 KB boot ROM, 64 KB SRAM, a 16550-compatible UART, a Core Local Interruptor (CLINT) with machine-mode timer, and the PNM router engine---all on a single die.

The choice of an SoC over a microcontroller is deliberate and driven by four requirements that no MCU satisfies simultaneously. First, OS compatibility: running Linux or Redox on the router chip---whether for development tooling, containerized dispatch services, or future userland extensions---requires either an Sv32 MMU or NOMMU kernel support; commodity MCUs typically provide neither, while RV32IMAFC cores with NOMMU Linux are a proven configuration\cite{29}. Second, host interface bandwidth: terminating a PCIe Gen5 endpoint requires dedicated PHY and controller blocks integrated into SoC-class silicon, not bit-banged GPIO on an MCU. Third, dispatch throughput: issuing 270 fabric dispatches per generated token at the sustained dispatch rates of Section \hyperref[sec:2.8]{2.8} demands clock rates and memory hierarchies beyond typical MCU envelopes. Fourth, routing-table capacity: source-routed tables for 512 nodes plus MoE expert maps exceed the on-chip SRAM of most MCUs but fit comfortably in SoC-class SRAM with cache backing.

The current RTL implements the RV32IMA baseline as a proof of concept; production silicon upgrades to RV32IMAFC (adding compressed instructions and single/double-precision floating point) to meet the NOMMU Linux requirement. The C firmware port targets both deployment modes: bare-metal static allocation for deterministic boot-time behavior, and a NOMMU Linux userspace daemon for operational flexibility.

\subsubsection*{2.17.1 Main thread placement: the router chip hosts transpiler orchestration}\label{sec:2.17.1}

For source-language transpilation (Section \hyperref[sec:2.7]{2.7}), the main thread runs on the router chip's CPU, not on any node PE or the host. This placement follows from three constraints. First, the compiler's IR lowering stage needs visibility into the full chassis topology---which nodes hold which compute units, which routing patterns connect them---and only the router chip holds the complete POST-discovered topology map. Second, the transpiler's output is a sequence of .pnm directives that must be injected through the fabric's ingress port, which terminates at the router chip's PCIe interface. Third, the doorbell-driven execution model requires a coordinator that tracks per-node kernel completion; placing this coordinator on a PE tile would consume a scarce compute unit for bookkeeping.

Concretely: the host compiles R/Haskell/HLSL source to typed IR (the existing CompileR/CompileHaskell/CompileHLSL pipeline), ships the IR to the router chip over PCIe, and the router chip's firmware lowers it onto the chassis---allocating registers, assigning kernels to nodes, computing routes, and emitting flits through its own PNM register block. The host never touches the spine directly. This keeps the PCIe link carrying compact IR rather than raw .pnm text, and it means the router chip can re-plan mid-execution if a node fails (Section \hyperref[sec:4.1]{4.1}), without round-tripping through the host.

The CPU core is a 5-stage finite-state-machine design (Fetch, Decode, Execute, Memory, Writeback) supporting the RV32I base integer instruction set, the M extension for single-cycle multiplication, and the Zicsr extension for machine-mode control and status registers. The core implements ECALL, EBREAK, and MRET for trap handling, and provides a combinational fetch address output for zero-cycle ROM access. Multi-cycle memory operations use a two-phase handshake: cycle 1 presents the address and write data, cycle 2+ waits for bus-ready, ensuring correct timing across all peripheral types. The ALU is fully combinational, driven by a dedicated input multiplexer that selects operands based on the decoded instruction, avoiding mixing blocking and non-blocking assignments in the sequential datapath.

The memory map partitions the 32-bit address space into five regions: Boot ROM at 0x0000\_0000 (64 KB, combinational read for zero-penalty fetch), UART at 0x1000\_0000 (16550-compatible, 8N1, configurable baud), CLINT at 0x2000\_0000 (machine-mode timer with software interrupt), SRAM at 0x8000\_0000 (64 KB, byte-enable write), and PNM router registers at 0xF000\_0000 (control, layer/module/length/data registers, status, error counters, dispatch and weight-flit counters). All peripherals use a valid/ready handshake on a shared bus; slave selects are decoded combinationally from the bus address, and read data is multiplexed to the CPU.

The RTL ships four SoC tiers sharing the CPU and bus fabric, so a deployment scales the management plane without changing firmware structure above the register layer. The BMC tier (bmc\_orchestrator\_top.v) is the five-region minimal map above. The MCU tier (orchestrator\_mcu.v) drops the CLINT for pure bare-metal dispatch loops. The MoE SBC tier (orchestrator\_sbc\_moe.v) adds a 64 KB BF16 gating SRAM and the moe\_gating unit so expert selection runs in silicon next to the dispatcher. The general SBC tier (orchestrator\_sbc.v) is the full machine: SRAM at 0x4000\_0000 sized by one synthesis-time parameter over a 500 KB--32 MB range (bare-metal dispatch loops at the bottom of the range; trimmed Linux tinyconfig or seL4 images at 1--4 MB; Linux plus CPython userland at 16 MB), a 1 GB LPDDR5 DRAM window at 0x8000\_0000 whose behavioral controller models row-activate plus CAS latency as a programmable parameter, a PCIe Gen5 x16 endpoint at 0xC000\_0000, an NVMe storage controller at 0xD000\_0000, and the PNM window at 0xF000\_0000. Each tier passes its own self-checking testbench (tb\_bmc\_orchestrator.v, tb\_orchestrator\_mcu.v, tb\_orchestrator\_sbc\_moe.v, tb\_orchestrator\_sbc.v).

Two of the SBC-tier blocks deserve specification because they change what the router chip can do unaided. The PCIe Gen5 x16 PHY (pcie\_phy.v) terminates the host uplink of Section \hyperref[sec:4.2]{4.2} in silicon rather than as a paper abstraction: an LTSSM walks Detect through L0, posted/completion headers stage through a completion FIFO whose pop is deferred one cycle to keep the ready/valid contract under spine backpressure (a ready/valid bug class the harness's audit log documents fixing), and a loopback parameter drives link bring-up without a live host---verified by twelve self-checking scenarios in tb\_pcie\_phy.v. The NVMe storage controller (nvme\_ctrl.v) attaches commodity M.2/U.2 SSDs behind a 64-deep command queue and 512-byte block engine, giving the router chip a persistence tier the fabric nodes deliberately lack: boot images and rootfs survive power cycles without re-streaming from the host, KV-cache offload spills to local flash instead of crossing the PCIe uplink, and the C firmware's block driver (pnm\_nvme.c) plus Lustre striping client (pnm\_lustre.c) speak the same registers from userland or bare metal. The result is a management plane that boots standalone---POST, routing tables, weights---while every byte of compute state remains in the chassis's commodity DRAM pools where Sections \hyperref[sec:2.4]{2.4} and \hyperref[sec:2.16]{2.16} place it.

The PNM router engine is a simplified version of the orchestrator\_chip.v fabric controller, driven by CPU register writes instead of a host PCIe stream. It implements a flit builder with CRC-16/CCITT-FALSE integrity checking, matching the fabric's wire format (Section \hyperref[sec:2.9]{2.9}). The CPU firmware runs the five-phase boot protocol of Section \hyperref[sec:2.5]{2.5} verbatim; a dedicated boot\_done detection circuit in the register block asserts when the firmware writes bit 2 of the ROUTE\_CTRL register, transitioning the chassis from initialization to the dispatch loop.

The design is verified by a self-checking testbench (tb\_bmc\_orchestrator.v) that loads a test program into ROM, exercises UART output, PNM register writes, and the boot-done signal, and checks for correct execution of LUI, ADDI, SW, and JAL instructions. The testbench instantiates the full SoC with a 2-layer, $2\times2$-node fabric configuration, stubbed spine injection/extraction, and all peripheral models.

Host access path. The management plane is shared between the router's own CPU and an external host through a two-port round-robin arbiter (pnm\_arb): master 0 is the on-chip RV32 core, master 1 is a SPI slave bridge (pi\_bridge) that lets a commodity single-board computer act as the chassis console. The bridge translates 48-bit mode-0 SPI frames---a header byte carrying read/write flag and register selector, a 32-bit data field, and a status acknowledgment byte---into valid/ready bus transactions; selectors 0--9 map to the standard PNM registers, and higher selectors pass through as raw byte offsets into an extended window holding a DMA engine: target address at 0x40, transfer word count at 0x44, control (start-write, start-read, done) at 0x48, and a FIFO data port at 0x4C. Both masters observe the same discipline---valid held until ready pulses---so either combinational or registered slaves can sit behind the arbiter without protocol change. The path is verified end-to-end by tb\_host\_bmc.v: SPI register writes land in the shared register file, a flit assembled by the bridge emerges from the spine ingress port with correct layer/module/CRC framing, boot-done asserts after the firmware sequence, and the dispatch counter reads back over the same serial link.

Memory-side verification. The DMA engine extends the identical discipline to node-local DRAM through the behavioral LPCAMM2/LPDDR6 module model (lpddr6\_camm.v), which implements the controller contract of Section \hyperref[sec:2.16]{2.16}---CAS-latency waits, one response per transaction, refresh bursts that stall command acceptance under the valid/ready handshake described there. tb\_dma\_lpddr6.v drives a four-word payload from the SPI host through bridge, arbiter, and FIFO into the module, verifies the array contents at the word-mapped address, counts exactly four writes at the module's activity counters, and reads the payload back byte-exact while the refresh scheduler absorbs more than two dozen stall events mid-transfer. Because the module is parameterized rather than hard-coded, the same testbench structure instantiates the same protocol stack against DDR4-class and DDR5-class SODIMM timing profiles (Section \hyperref[sec:3.7]{3.7}); all three generations pass identically, which is the machine-checked form of the claim that the chassis software stack is memory-generation agnostic. Where the deskside tier needs the controller itself rather than its timing alone, sodimm\_ctrl.v implements a real banked DDR controller behind the identical bus contract: a per-bank open-row tracker drives activate/precharge/CAS micro-sequences under enforced tRAS, so a row hit costs CAS latency, a miss adds tRCD, a conflict adds tRP+tRCD, and refresh closes every bank while never dropping an accepted command; tb\_sodimm\_ctrl.v checks those three latencies in closed form across DDR4-class and DDR5-class profiles, along with byte-enable masking and refresh-storm integrity.

\subsection*{2.18 PCB interconnect models}\label{sec:2.18}

The board-level wiring that Sections \hyperref[sec:2.3]{2.3} and \hyperref[sec:2.4]{2.4} describe qualitatively---spine mezzanine jumps, motherboard point-to-point traces, CAMM2 socket escapes---is modeled in the repository as parametric timed wires (pcb\_link.v). Each link is a shift-register delay line of $\lceil$DELAY\_NS/10$\rceil$ reference-clock cycles: data, valid, and the start/end-of-packet markers traverse bit-exactly, arriving a deterministic number of cycles later, and nothing else happens---no protocol conversion, no buffering, no reordering. Drive-current and trace-impedance parameters ride along as documentation of the electrical envelope each profile assumes. Three profiles cover the chassis's three PCB populations: the spine mezzanine connector ($\approx$5 mm mated, 128-bit), the motherboard point-to-point trace between repeater and node ($\approx$40 mm, 128-bit), and the CAMM2 socket escape to the MAC ASIC PHY ($\approx$20 mm, 192-bit). A wrapper (pcb\_triple\_link) instantiates three links sharing one clock domain for the X, Y, and spine egress fans a node actually has, with an aggregate flit counter for activity accounting.

tb\_pcb\_link.v pins the model down with four checks: a single-flit passthrough preserving SOP/EOP and payload; an exact-latency check that a 40 ns profile delivers in precisely its four configured stages, not approximately; a 200-packet randomized stream (1,408 bytes, per-packet lengths of 1--12 words) delivered byte-exact and in order against a scoreboard; and the triple-link aggregate counter matching an independent injection count. These bounds are what let Section \hyperref[sec:3.4]{3.4} add per-population wire delay into the closed-form latency expression as a compile-time constant rather than a statistical spread: each populated segment contributes exactly its configured delay, and the co-simulation's hop accounting already treats every such segment as one cycle of the closed form.

\subsection*{2.19 Inter-chassis optical links}\label{sec:2.19}

The scale-out path of Section \hyperref[sec:3.7]{3.7} leaves the chassis through the spine root, and its physical realization is a fiber-optic inter-chassis link pair (optical\_link.v). The link is deliberately protocol-transparent: it moves an uninterpreted byte stream with start/end-of-packet delimiters, so the fabric wire format of Section \hyperref[sec:2.9]{2.9} crosses the fiber untouched and neither endpoint rewrites a flit. Each unidirectional link is a dual-clock elastic buffer whose write and read pointers cross the domain boundary as Gray codes through two-flop synchronizers, followed by a propagation pipeline that models electrical-to-optical conversion, fiber transit at the speed of light in glass ($n \approx 1.47$, approximately 4.9 ns per meter), and optical-to-electrical recovery, quantized up to whole reference-clock cycles. A wrapper composes two links into one bidirectional chassis-to-chassis port---one fiber per direction, as an optical transceiver is actually wired.

Flow control across the fiber is structural rather than probabilistic. The read domain pops whenever the buffer is non-empty---no return handshake ever crosses the clock boundary---and all backpressure is applied on the write side, where transmit ready deasserts once the Gray-synchronized occupancy reaches buffer depth minus the propagation pipeline minus a guard term. Any producer that respects ready therefore cannot overflow, and a protocol violation latches a sticky error flag instead of silently dropping bytes. Idle-link latency measures propagation depth plus five cycles at the 100 MHz reference (synchronizer visibility plus the output register), which for the default two-meter board-to-board span budgets roughly 70 ns end to end---about one intra-chassis spine hop per meter of fiber, so a multi-chassis rack adds single-digit hop-equivalents to the closed-form latency of Section \hyperref[sec:3.4]{3.4}. Because an optical hop adds conversion and serialization error sources mid-fabric, the receiving router re-runs the CRC-16 check of Section \hyperref[sec:2.9]{2.9} on arrival (the shared crc16 engine) and refuses corrupt frames exactly as a node doorbell would, keeping silent corruption from propagating into a second chassis. A six-scenario testbench (tb\_optical\_link.v) pins the model down: exact idle-link latency, 200-word burst framing alignment against a scoreboard, a 600-word back-to-back flood delivered with zero loss, sparse packets separated by idle gaps and held in order, an asynchronous 100 MHz-to-80 MHz clock-domain crossing under sustained overload that observes transmit ready deassert, and simultaneous bidirectional traffic across the pair.

\section*{Results}\label{sec:results}

Cost, power, and thermal figures in this section are first-order engineering estimates intended to establish plausibility and sizing rules, not measured results; no physical prototype exists yet. The transport-level claims are machine-checked, not estimated: Section \hyperref[sec:3.5]{3.5} reports the co-simulation verification results against the cycle-exact fabric model of Section \hyperref[sec:2.15]{2.15}.

\subsection*{3.1 Reference chassis}\label{sec:3.1}

Table 2 summarizes the reference configuration: 512 nodes, 64 TB aggregate capacity, approximately 131 TB/s aggregate node-local bandwidth, approximately 197 TFLOPS FP64, and 8--10 kW chassis power, with a spine provisioned at $\approx$2 TB/s per direction from the cross-layer sizing rule of Section \hyperref[sec:2.1]{2.1}. The configuration is defined by three physical properties that no market correction can alter: (a) no interposer dependency---each node is a standard JEDEC CAMM2 module paired with a socketed DUV ASIC, with no advanced packaging between them; (b) low thermal density---approximately 15 W per node versus 300 W for an H100-class GPU, enabling air or liquid cooling without exotic thermal solutions; and (c) mature-node fabrication---the 14 nm DUV MAC ASIC benefits from yields above 90\% on established process lines, while the HBM stacks it pairs with consume roughly three times the wafer area per bit of commodity DRAM and absorb a disproportionate share of leading-edge fab capacity\cite{25}. Cost is a consequence of these physical properties: commodity DRAM in volume, mature-node silicon, and standard packaging yield a node at approximately \$2,000, but the structural advantage exists independently of DRAM spot pricing.

Table 2. Reference chassis configuration (first-order estimates).

Quantity | Reference value
Node | 128 GB LPDDR6 CAMM2 + 14 nm DUV MAC ASIC (192 FP64 FMA @ 1 GHz $\approx$384 GFLOPS FP64; $\approx$768 GFLOPS FP16; $\approx$1.5 TOPS INT8)
Node memory bandwidth | $\approx$256 GB/s (192-bit LPDDR6 CAMM2 @ 10.7 Gb/s/pin)
DDR5 baseline (contrast) | 192 GB dual-slot (2$\times$96 GB DDR5-5600 SODIMM); $\approx$90 GB/s per node---$\approx$3$\times$ below the LPDDR6 reference
Node fabric link | $\approx$16 GB/s (128-bit source-synchronous @ 1 GHz)
Node power | $\approx$15 W
Board (1 layer) | 8$\times$8 = 64 nodes; 8 TB; $\approx$16 TB/s; $\approx$1 kW
Chassis (8 layers) | 512 nodes; 64 TB; $\approx$131 TB/s; $\approx$197 TFLOPS FP64; $\approx$8--10 kW
Spine | Parallel 16 GB/s lanes; $\approx$2 TB/s per direction aggregate (cross-layer sizing rule, Section \hyperref[sec:2.1]{2.1})

The DDR5 baseline row contrasts the reference node with the commodity SODIMM deployment of Section \hyperref[sec:3.7]{3.7}: the same 14 nm ASIC, repeater, doorbell, and firmware against a dual-slot DDR5-5600 SODIMM pair ($\approx$90 GB/s, 192 GB). Section \hyperref[sec:3.2]{3.2}'s roofline places both flagship workloads below even that reduced ridge point---weight streaming at 1--2 FLOP/byte is bandwidth-bound on DDR5 exactly as on LPDDR6---so the bandwidth-bound conclusion is generation-independent, not an artifact of the LPDDR6 choice. The reference chassis simply selects the highest-bandwidth socketable module class for the same silicon.

\subsubsection*{3.1.1 The reference chassis as a dedicated capacity-serving appliance}\label{sec:3.1.1}

At full scale the 512-node, 64 TB chassis is positioned strictly as a serving appliance for capacity-bound, weight-stationary workloads---Mixture-of-Experts transformer inference chief among them\cite{27}---and the placement tooling makes that positioning concrete rather than aspirational. Compiling an MoE transformer onto the 8$\times$8$\times$8 chassis emits three artifacts: a routing table of exactly 512 entries---one 11-bit bitmap per repeater and node, loaded during POST; an expert map placing all experts of every model layer across the expert-parallel node pool with no data-parallel replication waste; and a per-token dispatch plan. For the Gemma-4-class reference model fixture (30 layers, 128 experts, top-8), the plan is 270 dispatches per token: 30 dense attention projections plus 240 expert activations, each carrying its destination coordinate, compute-unit assignment, KV-cache disposition, and flit byte count. The firmware boot sequence is verified against the same artifacts: five phases (POST discovery finding all 512 nodes, routing-table load of 512 bitmaps, weight upload of 8,045 integrity-checked commands, MoE gating load of 3,840 expert mappings---128 experts $\times$ 30 layers---and ready) complete with zero error counters, and the dispatch plan verifies coverage: every activated expert maps to a resident node within budget.

The capacity arithmetic is the point of the configuration. Each node reserves up to 80\% of its post-weight budget for KV offload, capped at 16K entries per resident layer so that context growth evicts cleanly instead of fragmenting; the remaining capacity holds expert weights at BF16. Because experts are placed round-robin-free---the map assigns each expert to exactly one node---a token's top-8 fan-out touches eight nodes plus one dense node per layer, and the closed-form route of Section \hyperref[sec:2.8]{2.8} bounds each hop regardless of which experts the gate selects. Nothing in the configuration is specific to this model family: the compiler re-derives placements, bitmaps, and dispatch plans from any config.json/safetensors pair, which is what makes the appliance a configuration of the architecture rather than a distinct machine.

\subsection*{3.2 Roofline: why mature nodes suffice}\label{sec:3.2}

The reference node's machine balance is precision-dependent, because the DUV MAC array packs INT8 and FP16 lanes at 4$\times$ and 2$\times$ the FP64 rate:
$$
B_{\mathrm{machine}}^{(\mathrm{FP64})} = \frac{384\ \mathrm{GFLOP/s}}{256\ \mathrm{GB/s}} \approx 1.5\ \mathrm{FLOP/byte}, \qquad
B_{\mathrm{machine}}^{(\mathrm{FP16})} \approx 3.0, \qquad
B_{\mathrm{machine}}^{(\mathrm{INT8})} \approx 6.0.
$$
A 7-point FP64 stencil has arithmetic intensity $\approx$0.3--0.5 FLOP/byte, below the 1.5 FP64 balance (strongly bandwidth-bound)\cite{5}. MoE expert weight-streaming sits near 1--2 FLOP/byte per token---below the 3.0 FP16 balance and far below the 6.0 INT8 balance, so MoE is bandwidth-bound on the precision it actually runs, even at the top of its intensity range. Performance is therefore limited by the LPDDR6 bus, never by the MAC array, for both flagship workloads. This is the load-bearing argument for DUV: a 14 nm or 28 nm array already saturates the memory it is attached to.

\subsection*{3.3 Cost and capacity model}\label{sec:3.3}

The cost advantage is structural, not cyclical. A node bill of materials: LPDDR6 CAMM2 module $\approx$\$1,900 ($\approx$\$15/GB at 2025--26 DRAM spot pricing), 14 nm MAC ASIC $\approx$\$40, socket/board/assembly share $\approx$\$60, for a node total of $\approx$\$2,000 ($\approx$\$16/GB). A 512-node chassis is $\approx$\$1M for 64 TB. The comparison to HBM is capacity-equivalent, not bandwidth-equivalent: a 64 TB HBM system requires 819 H100-class GPUs (80 GB each)---$\approx$\$24.6M of accelerators at $\approx$\$30k apiece, $\approx$0.25 MW of GPU power alone---against $\approx$\$1M, $\approx$10 kW, and one chassis for the same 64 TB here. Even the HBM itself, at $\approx$\$375/GB\cite{2}, costs $\approx$\$24M for 64 TB: 24$\times$ the entire PNM chassis, memory included. The HBM fleet's aggregate bandwidth dwarfs the 131 TB/s of node-local LPDDR6, but Section \hyperref[sec:4.5]{4.5} shows inference serving is capacity-bound rather than bandwidth-bound, so that excess bandwidth idles during decode while the capacity bill is paid in full. DRAM is a cyclical commodity\cite{25}, and the 2025--26 spot prices used here reflect a shortage-driven spike; but the structural factors---commodity modules versus interposer-bound stacks, mature-node DUV versus EUV, standard packaging versus advanced integration---persist across market cycles. The advantage is architectural, not a bet on a particular price point.

\subsection*{3.4 Thermal and latency budgets}\label{sec:3.4}

Thermal: $\approx$15 W per node, $\approx$1 kW per board, $\approx$8--10 kW per chassis; 15 L/min of 30$^\circ$C water with 10$^\circ$C rise and orifice-balanced branches keeps worst-case DRAM junction $\leq$70$^\circ$C against an 85$^\circ$C throttle.

Latency (bounded, not average): lxy\_repeater compare $\approx$0.1 ns; HFR hop $\approx$1--2 ns; worst-case on-board route (8$\times$8) 14 hops $\approx$14--28 ns; spine traversal (7 layer hops) $\approx$7--14 ns; link serialization for a 4 KB message at the 16 GB/s fabric link $\approx$256 ns (the dominant term); doorbell wake $\approx$1 ns. The end-to-end activation bound for a 4 KB cross-node message is therefore $\approx$0.3 $\mu$s, dominated by link serialization, and it scales linearly with message length ($\approx$64 ns per KB). The closed form verified in Section \hyperref[sec:4.3]{4.3} is the sum of these terms: hop count $\times$ per-hop delay plus message bytes divided by the link rate; the co-simulation checks the hop-count and serialization terms exactly and treats contention as a separate, compile-time-scheduled term.

\subsection*{3.5 Co-simulation verification results}\label{sec:3.5}

Table 3 summarizes the co-simulation runs on the 48-node chassis (3$\times$4$\times$4, deterministic seed, three parallel fabric slices). Every scenario passes all five oracle checks of Section \hyperref[sec:2.15]{2.15}.

Table 3. Co-simulation verification results (48-node chassis, 3$\times$4$\times$4, seed 1, three parallel slices).

Scenario | Routed | Pass-through | Wire bytes | Activations | Corrupt refused | Latency min/mean/max (cycles)
sweep | 48 | 0 | 1,104 | 48 | 0 | 24/25/27
vcsweep | 48 | 24 | 1,656 | 48 | 0 | 24/25/27
load | 450 | 50 | 18,281 | 450 | 0 | 9/149/558
hotspot | 380 | 20 | 151,191 | 380 | 0 | 27/2,888/8,279
stress | 1,087 | 65 | 164,914 | 1,051 | 36 | 12/671/8,279
replay | 289 | 11 | 52,132 | 285 | 4 | 14/758/8,255

sweep's mean latency of 25 cycles matches the closed form of Section \hyperref[sec:3.4]{3.4} exactly, so per-hop accounting is exact rather than approximate. load delivers all 450 routed flits byte-exact despite sinks that accept only one byte in four; its 558-cycle worst case is a backpressure artifact of a slow consumer, not a routing anomaly. hotspot's 8,279-cycle worst case on its corner node is exactly the serialization and DMA-throttled drain of kilobyte-class tokens behind the 1-in-8 gate---and not a single byte is dropped. stress refuses all 36 corrupt-CRC flits---the count matching the software twin's refusals and the stub's corrupt\_out verdicts exactly---while the remaining 1,051 activations fire kernels whose results are byte-exact against the golden model. replay's two runs produce 2.1 MB of bit-identical delivery logs across three slices, machine-checking deterministic replay. vcsweep's routed flits arrive within the same 24--27 cycle window despite concurrent pass-through traffic across all four virtual-channel classes, confirming the arbitrated egress tree adds no measurable delay at this load.

The 512-node chassis (8$\times$8$\times$8, eight parallel slices) exercises the full spine at scale under all six scenarios of Section \hyperref[sec:2.15]{2.15}: sweep at 24/27/31 cycles (min/mean/max), vcsweep likewise at 24/27/31, load with 18,782 wire bytes at 18/157/583 cycles, hotspot at 25/2,878/8,311 cycles, stress refusing 79 corrupt messages with 2,773 byte-exact activations at 12/632/8,287 cycles, and replay reproducing all eight slices' delivery logs bit-identically across two runs (2,053,924 bytes). The compiler path adds a two-expert program---a byte-sum kernel with silicon bias +7 and a 16-weight dot kernel with bias +3---on the same chassis: five dispatched tokens produce five doorbell activations, five byte-exact kernel results, zero refusals, and 14--26 cycle latencies.

The reasoning behind these checks is what makes them sound. Byte conservation is only testable because routing is single-path and deterministic: on a fabric with multiple legal routes the oracle would need a routing model of its own, whereas here the injected byte count is the delivered byte count by construction, and any divergence is a gate-level bug. The doorbell accounting is sound because the hardware verdict and the software twin are independently implemented---the HDL doorbell and CRC datapath on one side, the Go twin on the other---so their agreement is evidence, not tautology. Replay is sound because the fabric has no hidden state: if any queue, counter, or configuration bit persisted between runs, the second delivery log would diverge. And the sweep latency check is exact because contention-free schedules are compile-time properties (Section \hyperref[sec:2.8]{2.8}): the simulator runs the same schedule the AOT compiler emits, so the closed form is checked against the gates, not against a second software model. The same argument covers the arbitrated reverse path: the DES of Section \hyperref[sec:2.15]{2.15} shares no code with the Verilog merge tree yet reproduces its delivery stream byte-for-byte on a mixed four-class program and matches the same closed form on an idle chassis, so the egress arbitration is pinned down by two independent implementations.

The host-to-chassis path adds four self-checking testbenches outside the co-simulation harness, each pinning down one interface boundary of Sections \hyperref[sec:2.17]{2.17}--2.18. tb\_host\_bmc.v drives the SPI console through bridge and arbiter into the shared register file and observes a correctly framed flit leave the spine ingress, boot-done assert, and the dispatch counter read back over the same serial link. tb\_dma\_lpddr6.v moves a four-word payload host$\to$LPDDR6$\to$host through the extended register window: the module's activity counters report exactly four writes and four reads, array contents match at the word-mapped address, and the readback is byte-exact across 27 absorbed refresh stalls. tb\_sodimm\_lpddr.v repeats the identical sequence against DDR4- and DDR5-class SODIMM timing profiles (CAS waits of 15 and 36 cycles against LPDDR6's 4, refresh intervals tightened to 780 cycles) and passes all three generations unchanged---the protocol stack carries no timing assumption of any single DRAM generation. tb\_pcb\_link.v verifies the interconnect models of Section \hyperref[sec:2.18]{2.18}: exact multi-stage latency, a 200-packet/1,408-byte stream delivered byte-exact in order, and triple-link activity counters matching injection counts.

\subsection*{3.6 The DUV MAC ASIC: detailed specifications}\label{sec:3.6}

Table 4 is the complete datasheet outline for the reference node's ASIC (Variant A of Section \hyperref[sec:2.4]{2.4}). Every figure is a first-order estimate consistent with the node budget of Table 2.

Table 4. DUV MAC ASIC specifications (reference node; first-order estimates).

Parameter | Specification
Process | 14 nm DUV (28 nm option), mature-node, high-yield; $\approx$60 mm$^2$ die---comfortably inside the DUV reticle and the mature-node yield sweet spot
Clocks | 1 GHz fabric/MAC clock; 10.7 Gb/s/pin LPDDR6 I/O via on-die PLL; source-synchronous fabric forwarding
MAC array | 192 FP64 FMA lanes; 384 FP16 lanes (2$\times$ packing); 768 INT8 lanes (4$\times$ packing)
Peak throughput | $\approx$384 GFLOPS FP64; $\approx$768 GFLOPS FP16; $\approx$1.5 TOPS INT8
Memory interface | 192-bit LPDDR6 controller, 4$\times$48-bit channel groups, 10.7 Gb/s/pin, $\approx$256 GB/s; bus SECDED plus on-die ECC pass-through (Sections \hyperref[sec:2.10]{2.10} and \hyperref[sec:2.16]{2.16})
Fabric interface | AXI-Stream slave (RX) and master (TX), 128-bit @ 1 GHz, $\approx$16 GB/s per direction; wormhole framing, link-local CRC per flit
Doorbell DMA | byte counter, running CRC, DEST compare; DOORBELL\_TRIG/DOORBELL\_ACK and NODE\_ERR on the sideband (Section \hyperref[sec:2.9]{2.9})
On-chip SRAM | $\approx$1--2 MB: elastic input buffer of at least one maximum-length message (compiler-set capacity), MAC staging, CRC/DMA workspace
Power | $\approx$10 W total ($\approx$9 W MAC array, $\approx$1 W fabric I/O) at 0.8 V core, 1.2 V I/O, 1.1 V LPDDR6 PHY
Package | Variant A: socketed LGA/card-edge, $\approx$800--850 pins, $\approx$20 mm from the CAMM2 socket; Variant B: mounted on the module substrate

The pin groups of Table 4 map onto the socket as follows: LPDDR6 PHY $\approx$290 pins (192 DQ, 24 DQS pairs, CA/CS/CK); fabric $\approx$270 pins (128-bit data in each direction plus valid/sop/eop/ready and a forwarded clock); sideband 4 pins; clock and test (reference oscillator, JTAG, configuration straps) $\approx$20; the remaining $\approx$250 pins are power and ground---the majority of any high-speed package. All high-speed groups are length-matched within the package substrate, and the PHY-to-socket escape is a controlled-impedance bundle consistent with the board outline of Section \hyperref[sec:2.4]{2.4}.

Two sizing remarks complete the picture. First, the MAC array is the die's minority: at 14 nm, 192 FP64 FMA lanes plus the 4$\times$ INT8 packing occupy roughly a third of the $\approx$60 mm$^2$, the LPDDR6 PHY and its escape another third, and SRAM, fabric I/O, and clocking the remainder---so the silicon is memory-interface-bound, exactly as the roofline of Section \hyperref[sec:3.2]{3.2} predicts, and there is ample DUV margin to grow the array without moving the ridge point. Second, the ASIC holds no instruction path and no operating-system interface; its only program is the routing-table entry, resident kernel weights, and the KERNEL\_CONST bias constant loaded through the POST sequence of Section \hyperref[sec:2.5]{2.5}, which is what keeps the pinout static across every workload the chassis can run.

The central router chip is a second, distinct die: a RISC-V BMC/SoC that owns the chassis management plane (Section \hyperref[sec:2.17]{2.17}). Unlike the compute nodes, which are fixed-function MAC ASICs, the router chip is fully programmable---it runs firmware from on-chip ROM, handles PCIe Gen5 x16 host uplinks through its integrated PHY (Section \hyperref[sec:4.2]{4.2}), performs POST discovery, evaluates MoE gating networks, attaches NVMe storage for its own persistence tier, and dispatches tokens into the fabric. The router die is smaller than the compute die (no MAC array, no 192-bit LPDDR6 PHY; its DRAM interface is a single LPDDR5-class controller window sized in hundreds of megabytes), uses the same mature-node process, carries an SRAM bank synthesized anywhere in a 500 KB--32 MB range by board tier, and is socketed and cold-swappable for serviceability.

The MAC array is not a monolithic unit but a collection of specialized compute units, each optimized for a different precision and operation class. The reference node instantiates eleven unit types, selected at compile time by the kernel mapper (Section \hyperref[sec:2.7]{2.7}):

Table 4a. Compute unit taxonomy (per-node; reference configuration).

Unit | Precision | Pipeline latency | Operation class | Workload mapping
fp16\_fma | FP16 | 3 cycles | Fused Multiply-Add | FP16 transformer inference, mixed-precision stencil halos
bf16\_fma | BF16 | 3 cycles | Fused Multiply-Add | MoE expert weight-streaming, dense MLP projections (default)
fp32\_fma | FP32 | 3 cycles | Fused Multiply-Add | High-precision layer norms, attention softmax denominators
fp64\_fma | FP64 | 3 cycles | Fused Multiply-Add | FP64 stencil computation, scientific HPC kernels
fp16\_mac\_array | FP16 | variable | Systolic MAC array | FP16 attention QKV projections
bf16\_mac\_array | BF16 | variable | Systolic MAC array | BF16 attention QKV projections (default)
fp32\_alu | FP32 | 3--28 cycles | ALU (ADD/SUB/MUL/DIV/MIN/MAX/CMP) | Layer norm, softmax, activation functions
int8\_mac | INT8 | 2 cycles | Multiply-Accumulate | Quantized inference, INT8 weight-stationary kernels
int4\_mac\_array | INT4/INT8 | variable | Systolic MAC array | INT4 quantized inference (4x density vs BF16)
fp4\_mac\_array | FP4 (E2M1) | variable | Systolic MAC array | FP4 quantized inference (4x density vs BF16)
mxfp4\_mac\_array | MXFP4 | variable | Systolic MAC array | MXFP4 block-scaled (OCP microscaling) inference

The `fp32\_alu` unit deserves particular attention: it provides the division operation that the FMA units lack, and division is essential for normalization layers (layer norm, batch norm, softmax denominator). The ALU implements restoring binary long division with a 28-cycle latency, alongside 3-cycle MIN/MAX/CMP and 4-cycle FMA-path ADD/SUB/MUL. A dedicated `fp32\_alu\_chip` wrapper integrates the ALU with the PNM fabric's AXI-Stream interface: it deserializes 4-byte FP32 operands from the flit byte stream, feeds them to the ALU, and serializes the result back, all while running the CRC half of the doorbell discipline in hardware. This chip is the compute unit that sits behind the PE tile stub (Section \hyperref[sec:2.9]{2.9}) for layer norm and activation workloads.

The quantized precision spectrum extends the taxonomy down to 4-bit floating point. The `fp4\_mac\_array` and `mxfp4\_mac\_array` units operate on FP4 (E2M1) values packed two per byte, dequantizing each nibble to fixed point and accumulating to INT32---4x the packing density of BF16, equal to that of INT4, at the cost of reduced input precision. Following the OCP Microscaling convention, MXFP4 additionally carries one shared 8-bit block scale per 32-element block, so a run of weights can recover the exponent range that a bare 4-bit field cannot represent. A 4-bit value is too narrow to multiply directly, so both units follow the same dequantize-then-multiply convention as the INT4 path: the systolic array accumulates fixed-point products in INT32 and the wider accumulator absorbs the reduced input precision, with the block scale folded into the mantissa shift. The `weight\_dequant` unit was extended to the same four modes (INT4, INT8, FP4, MXFP4), and the co-simulation `QuantMode` offers `int4`, `fp4`, and `mxfp4` weight-only options alongside the default BF16. These narrower formats are a deployment option for bandwidth-bound weight streaming, not the default: BF16 remains the reference precision, and the closed-form latency bounds of Section \hyperref[sec:3.4]{3.4} hold for every unit type because the array pipeline depth, not the data width, sets the critical path.

All FMA modules share an identical interface: `clk, rst\_n, a, b, c, valid\_in` producing `result, valid\_out` with a 3-cycle pipeline. The systolic MAC arrays have a variable latency depending on the matrix dimensions. The PE tile stub (Section \hyperref[sec:2.9]{2.9}) uses the `USE\_FMA` parameter to select between bias-add mode (element-wise addition of a compile-time constant) and BF16 FMA mode (weight-stationary multiply against a resident weight). The model compiler (Section \hyperref[sec:2.7]{2.7}) assigns unit types based on the operation: MoE experts use BF16 FMA (weight-stationary), attention projections use BF16 systolic arrays, layer norms use FP32 ALU (for division), and dense MLP layers use BF16 FMA. This assignment is static and immutable: the same node always runs the same unit type for the same workload, which is what makes the deterministic latency bounds of Section \hyperref[sec:3.4]{3.4} possible.

\subsection*{3.7 Configuration spectrum: SODIMM deskside units to spatial supercomputers}\label{sec:3.7}

The architecture is a family of machines, not a single SKU, and the degrees of freedom are deliberately orthogonal. Four axes vary independently without touching the others: the memory module generation on each node, the compute-unit taxonomy instantiated behind the PE tile, the router-chip tier managing the chassis, and the physical extent (boards stacked into layers). Every axis is exercised by running RTL or tooling in the repository, and every combination speaks the same wire format (Section \hyperref[sec:2.9]{2.9}), honors the same doorbell discipline, and compiles from the same .pnm program schema---which is the concrete meaning of modularity here.

Memory axis: CAMM2 versus SODIMM. The flagship node pairs the MAC ASIC with LPCAMM2/LPDDR6\cite{16,17} for the $\approx$256 GB/s of Table 2, but the identical ASIC, repeater, doorbell, and firmware run against commodity SODIMM modules\cite{30} where a deployment calls for desktop economics instead of density. A dual-channel DDR5-5600 SODIMM provides $\approx$45 GB/s per module with capacities to 96 GB today; a dual-slot node reaches $\approx$90 GB/s and 192 GB. Node bandwidth drops roughly $3\times$, but Section \hyperref[sec:3.2]{3.2}'s roofline already places both flagship workloads below even the reduced ridge points: weight-streaming at 1--2 FLOP/byte remains bandwidth-bound on DDR5 exactly as on LPDDR6, under the same static schedule and the same closed form. The claim is machine-checked rather than asserted: tb\_sodimm\_lpddr.v drives one unmodified protocol stack against DDR4-, DDR5-, and LPDDR6-class timing parameters and verifies byte-exact DMA in all three generations---what changes between them is only the module's wait-state count and refresh cadence, precisely the parameters the behavioral module exposes.

Scaling down: GPU-class deskside units. The spine is optional hardware below chassis scale. A single 4$\times$4 board---16 nodes---forms a deskside unit in the multi-TB class (up to 2 TB at CAMM2 densities; 1.5 TB on 96 GB SODIMMs) drawing under 250 W: the power envelope of one workstation, the capacity of a small rack, and no spine stack at all. With every destination on the sender's own layer there is nothing above the board for the vertical fabric to do: the latency closed form loses its spine term (zero effective spine hops), virtual-channel classes 1 and 2 go uninstantiated, and the spine may be omitted entirely---the board's egress merge trees terminate directly in the router instead of an lxy\_repeater column, or retain that one column as a passive turn-around stub if a later stack-up is anticipated. Nothing else changes: source-routed matching compares coordinate bytes against local parameters, never absolute topology position, so the wire format, doorbell discipline, and .pnm program schema are byte-for-byte those of Section \hyperref[sec:2.9]{2.9}. The MCU-tier router variant (Section \hyperref[sec:2.17]{2.17}) suffices at this scale---16 routing entries fit in 4 KB of SRAM---and a single-board computer console reaches the PNM window over the SPI bridge of Section \hyperref[sec:2.17]{2.17} without PCIe. Workloads whose resident set fits a few hundred GB---MoE expert pools, FP64 stencil subdomains, seismic image stacks---serve interactively from this unit; the compile-model/run-driver pipeline targets it by passing board dimensions, with no source change. At this end of the spectrum the machine competes not with datacenter HBM but with the local-GPU status quo quantified in Section \hyperref[sec:3.3]{3.3}: a workstation GPU holds at most 80--96 GB behind its 300 W package, while the deskside unit holds an order of magnitude more capacity for comparable power.

Scaling out: spatial supercomputers. Chassis-level scaling replicates the reference configuration and joins replicas through a second routing level: each chassis keeps its deterministic internal closed form, and cross-chassis traffic leaves through the spine root as ordinary flits toward a chassis-interconnect tier---the extension point Section \hyperref[sec:4.4]{4.4} identifies. Because routing tables are immutable per boot session and schedules are compile-time artifacts, a multi-chassis job is assembled by concatenating per-chassis schemas at batch boundaries; no runtime coherence protocol appears anywhere in the stack. The result is a spatial supercomputer whose building block is a $\approx$\$1M, 10 kW board stack rather than an interposer-bound accelerator tray, so scale-out economics inherit the commodity-memory advantage rather than diluting it.

\section*{Discussion}\label{sec:discussion}

\subsection*{4.1 Failure Model and Recovery}\label{sec:4.1}

Losses are observable before they become failures. Every injected message has a compile-time latency budget, and its source expects DOORBELL\_ACK within that window (Section \hyperref[sec:2.4]{2.4}). A missing ACK---whether the message was dropped, truncated, or refused at the doorbell---raises NODE\_ERR for that destination; there is no silent hang. NODE\_ERR assertion triggers a four-step protocol: (1) quarantine the node from the live topology schema; (2) discard in-flight flits destined for it (the affected set is exactly enumerable on a single-path fabric); (3) recompute---the source re-injects the work toward a healthy replica or the reinitialized node, under the same doorbell discipline; this is functional recomputation, not bit-exact replay of the failed node's memory (Section \hyperref[sec:2.10]{2.10}); (4) re-export the amended schema at the next batch boundary. Routing tables are immutable within a boot session, so step (4) is a controller-coordinated reload, not runtime adaptation: the router regenerates the AOT tables against the amended schema and loads them during the batch boundary, when the fabric is quiet. The chassis continues in degraded mode with no restart. A failed spine lane is handled at the physical layer---periodic link re-training excludes the dead lane and the remaining lanes re-lane transparently---so routing tables need no change for lane degradation; total spine failure partitions the chassis (the Conclusion section).

\subsection*{4.2 Host Interface and I/O}\label{sec:4.2}

The central router chip terminates dual PCIe Gen5 x16-class uplinks ($\approx$128 GB/s aggregate) or 100/200 GbE equivalents. The host streams initial weights at boot, injects input tokens, and drains outputs. Runtime token traffic is kilobyte-class: the host injects one input token per generated token, and the 270 fabric dispatches per token (Section \hyperref[sec:2.6]{2.6}) are generated on-router from the resident expert map, so the uplink carries token-rate traffic---roughly $4\times10^{6}$ tokens/s of kilobyte-class tokens, about 4--30 GB/s---not dispatch-rate traffic, and remains well below the $\approx$128 GB/s uplink. The fabric's internal dispatch rate (Section \hyperref[sec:2.1]{2.1}) exceeds the uplink rate precisely because the per-token dispatch multiplier is realized inside the fabric. The host never touches the compute fabric directly.

\subsection*{4.3 Verification and Formal Properties}\label{sec:4.3}

Deadlock freedom (sketch). Every flit is assigned one of four virtual-channel classes in strictly increasing order: (0) on-board egress (X-then-Y dimension order), (1) spine ascent, (2) spine descent, (3) on-board delivery (X-then-Y from spine attachment to destination). Upward and downward spine flows are separate virtual channels (classes 1 and 2), which is why the count is four, not three. Intra-class dependencies are acyclic by standard results; cross-class dependencies only point from lower to higher classes. The global channel dependency graph is therefore acyclic, and by Dally and Seitz the fabric is deadlock-free\cite{13,14}.

The buffering that virtual channels require lives only at the two kinds of state points in the design: the lxy\_repeater spine attachments (per-class elastic buffers) and the node DMA ports (Section \hyperref[sec:2.9]{2.9}). HFRs remain stateless re-timers: they carry no per-VC state and never decide a route, which is what keeps the "no queues on the data path" claim true for the forwarding elements. The per-VC attachment buffers are credited in the flow-control sense below, so a blocked flit occupies its class's buffer and cannot block another class.

Livelock freedom follows from minimal single-path routes. Consumption is enforced structurally: the per-VC elastic buffers at spine attachments and node ports hold at least one maximum-length message, and sources may not inject without destination buffer credit, so no class ever exhausts and blocks a lower class. The egress merge tree adds no deadlock surface of its own: its dependency graph is a DAG---board egress toward spine ascent, never across spine descent---so the packet-atomic round-robin merges of Section \hyperref[sec:2.15]{2.15} cannot deadlock and the argument extends to the arbitrated reverse path. Worst-case latency for any coordinate pair is a closed-form expression---hop count $\times$ per-hop delay plus message length divided by link rate (Section \hyperref[sec:3.4]{3.4})---and contention is absent by construction because the AOT scheduler (Section \hyperref[sec:2.8]{2.8}) assigns message schedules that respect every buffer's credit budget; the co-simulation verifies the hop-count and serialization terms exactly and treats the scheduler's credit feasibility as a compile-time property. Given identical inputs onto an identical fabric, every computation and message schedule reproduces bit-exactly (Section \hyperref[sec:2.10]{2.10}).

\subsection*{4.4 Limitations and Future Work}\label{sec:4.4}

LPDDR6 (JESD209-6)\cite{17} CAMM2 parts and high-capacity modules are forward-looking; the current JEDEC socketable module class, LPCAMM2 (JESD318), is LPDDR5X and trims per-node bandwidth to $\approx$136--153 GB/s with 48--64 GB modules today\cite{16,17}. The architecture accommodates both module classes with no change to the MAC ASIC or routing fabric. Manifold balancing requires CFD and hardware validation. Multi-chassis scale-out needs a second routing level and one additional virtual-channel class; the physical interconnect itself---a verified protocol-transparent optical link pair with asynchronous elastic buffering (Section \hyperref[sec:2.19]{2.19})---is modeled in the repository, leaving the second-level router specification and the VC-field widening as the remaining design work. General lazy evaluation needs a dynamic heap this design omits. A dual-spine variant eliminates single-chassis partition mode. The liveness claims of Section \hyperref[sec:4.3]{4.3} remain an argument---a simulation can exercise the fabric but cannot prove acyclicity of its channel-dependency graph\cite{9}---while the transport claims (lossless delivery, byte-exact reproduction, doorbell refusal accounting) are machine-checked by the co-simulation of Sections \hyperref[sec:2.15]{2.15} and \hyperref[sec:3.5]{3.5}.

Validation scope. This paper's empirical claims rest on a validated behavioral substrate, not on fabricated silicon. The Verilog fabric, compute units, routers, and memory models are exercised by self-checking testbenches and the Go co-simulation at 100 MHz with one byte per cycle (Section \hyperref[sec:2.3]{2.3}); there is no synthesis, place-and-route, timing closure, or silicon measurement behind the bandwidth, power, thermal, or cost figures. The 1 GHz 128-bit links of Section \hyperref[sec:2.3]{2.3} are never clocked in simulation---the $\approx$160$\times$ gap is stated explicitly there---and the nanosecond latencies of Section \hyperref[sec:3.4]{3.4} are projections of the closed form at the characterized link rate. The DRAM, NVMe, and PHY blocks are behavioral stubs with parameterized CAS latency and refresh scheduling: they verify control flow, framing, and byte-enable behavior, not DRAM silicon. The MoE gating dispatch exercised by the co-simulation is a deterministic token-seeded hash, not a trained gating network (disclosed in Section \hyperref[sec:2.12]{2.12}); a learned 128-expert BF16 router would need roughly 0.3--1 MB of projector weights---5--15$\times$ the 64 KB gating SRAM---so the router sizing and gating evaluation are likewise projections. CRC-16/CCITT-FALSE is a probabilistic integrity check, not a proof: at the characterized fabric rate of $\approx$2.5$\times10^{8}$ 8 KB frames per second, the per-frame residual is 1 in $2^{16}$---approximately one undetected corruption every 50 days at a BER of $10^{-15}$, and one every few seconds at $10^{-12}$---so bit-exactness claims apply to the simulated traffic and to the residual rate stated here. The 512-node scale evidence is a routed wire record of 18 KB plus 2.0 MB across two replay runs; the co-simulation partitions the chassis into parallel slices, it does not instantiate 512 nodes of application state. Power figures are first-order design-point estimates consistent with the event-accounting checks of the repository's power monitor; the DRAM-read energy implied by full-bandwidth traffic is a projection to be validated, not a measurement, and thermal figures await the CFD validation of this section. BOM figures exclude mezzanine connectors, chassis assembly, and the isothermal manifold. What is machine-checked is exactly this: byte-exact delivery and replay, doorbell refusal accounting, and closed-form latency equality on the 48- and 512-node chassis (Section \hyperref[sec:3.5]{3.5}). The paper is therefore a position argument carried by a verified behavioral substrate, and its quantitative claims should be read as projections on top of that substrate.

Reliability claims and their validation scope. The only integrity mechanisms the co-simulation actually exercises are the end-to-end CRC-16/CCITT-FALSE check of Section \hyperref[sec:2.9]{2.9} and the link-local CRC of Section \hyperref[sec:2.10]{2.10}: the stress scenario injects corrupt-CRC flits and the doorbell refuses every one, with the hardware verdict matching the software twin's refusal count exactly (Section \hyperref[sec:3.5]{3.5}). LPDDR6 on-die ECC, the MAC ASIC's bus SECDED, and the background scrub pass of Sections \hyperref[sec:2.10]{2.10} and \hyperref[sec:2.16]{2.16} are design claims, not validated artifacts: the repository's DRAM and PHY models are behavioral stubs with no ECC datapath, no SECDED encoder or decoder exists in the RTL, and no testbench or co-simulation run exercises them---the scrub's mean-time-to-undetected-mis-correction bound is a sized projection, not a measured property. No functional-safety certification (IEC 61508 SIL, ISO 26262 ASIL, or equivalent) is claimed for the fabric or the router chip. The probabilistic residual of the frame CRC---1 in $2^{16}$ per frame at the rate stated in the previous paragraph---bounds every integrity claim in this section, and it applies unchanged to both memory baselines. A DDR5 baseline changes the ECC comparison but not its validation scope: commodity DDR5 RDIMMs provide conventional bus ECC with sideband check bits alongside the on-die refresh ECC that DDR5 carries internally; LPDDR6 relies on on-die ECC; and in both cases the controller SECDED of Section \hyperref[sec:2.10]{2.10} is the same unvalidated design claim.

DDR5 baseline economics. The SODIMM deployment of Section \hyperref[sec:3.7]{3.7} trades bandwidth for cost. A dual-slot DDR5 node carries $\approx$192 GB at $\approx$\$2/GB 2025--26 spot pricing (roughly \$400 of DRAM) against $\approx$\$1,900 for the 128 GB LPDDR6 CAMM2 of Table 2---a $\approx$7$\times$ lower memory cost per byte, partially offset by doubled module count and socket area at equal capacity. Node power stays inside the $\approx$15 W envelope: the DDR5 dual-channel pair draws roughly 5--8 W against $\approx$5 W for the LPDDR6 CAMM2 module. The bandwidth-bound conclusion of Sections \hyperref[sec:3.2]{3.2} and \hyperref[sec:3.7]{3.7} is unchanged, but the DDR5 node's $\approx$3$\times$ lower bandwidth raises the per-GB/s cost of the fabric and spine provisioned behind it, making the DDR5 baseline a capacity-maximizing rather than bandwidth-maximizing point on the spectrum of Section \hyperref[sec:3.7]{3.7}. Multi-chassis scale-out compounds the difference: the second routing level this section identifies is dimensioned from per-node bandwidth, so a DDR5 baseline carries $\approx$3$\times$ less cross-chassis bandwidth per chassis, and every published chassis figure (Table 2, Sections \hyperref[sec:3.5]{3.5} and \hyperref[sec:4.4]{4.4}) is stated for the LPDDR6 reference only---redeployment on DDR5 requires re-planning schedules, not re-architecting them, because the compile-time scheduler of Section \hyperref[sec:2.8]{2.8} absorbs the new bandwidth ratios unchanged. The 2025--26 DRAM spot prices behind both baselines are shortage-driven and cyclical (Section \hyperref[sec:3.3]{3.3}); the structural comparison---commodity modules and mature-node silicon against interposer-bound HBM---does not depend on them.

The central router chip is a single point of failure, and this paper treats that fact explicitly rather than designing around it. Chassis availability is bounded by the router's availability multiplied by the fabric's: the router is socketed and cold-swappable, POST re-discovers the chassis in under a second (Section \hyperref[sec:2.5]{2.5}), and all in-flight traffic drains through the stateless fabric while the spare boots, so a router replacement costs a sub-second dispatch outage. There is no live router failover---dual-router handover of the JIT tables is future work---and the reference design accepts the resulting availability bound in exchange for a single programmable silicon. An HFR, lxy\_repeater, or node failure degrades the machine by exactly that node (Section \hyperref[sec:4.1]{4.1}); a router failure halts dispatch until the spare is in place.

\subsection*{4.5 Impacts: capacity-bound workloads at commodity economics}\label{sec:4.5}

The workload landscape has bifurcated. On one side, training and interactive inference remain bandwidth-dominated, served by GPU accelerators with HBM. On the other side, a growing class of capacity-bound workloads---Mixture-of-Experts transformer inference, scientific HPC stencil computations, climate modeling, seismic imaging, and large-scale numerical simulation---are bottlenecked not by peak FLOPs but by the memory capacity required to hold resident data\cite{21,22,23,26}. These workloads share a common memory profile: the arithmetic intensity is low (0.3--2 FLOP/byte), the data sets are large (terabytes to petabytes), and the dominant cost is moving data from memory to compute, not the compute itself\cite{5}. NVIDIA already reported roughly 40\% of datacenter revenue from inference in 2024---'probably understated'---and its subsequent 'inference inflection' commentary is consistent with the broader shift toward capacity-bound serving\cite{2}.

Capacity-bound workloads have a memory profile that is the mirror image of training's. Training is bandwidth-dominated: gradients and activations must move at the peak rate the silicon can consume. Capacity-bound serving is weight-stationary: for every token the serving hardware streams the model's weights once from memory---decode arithmetic intensity is $\approx$2 FLOP/byte at batch 1 in FP8, far below the ridge point of any current GPU---and performs only a couple of operations per byte read\cite{24,26}. Serving throughput is bounded by memory bandwidth, and what can be served is bounded by memory capacity: fitting weights plus KV cache into HBM, not peak FLOPs, decides the batch a server can hold, which is why industry analyses call HBM capacity 'throughput capital'\cite{26}. MoE serving in particular is memory-capacity-bound: the binding resource is KV cache plus activated expert weights, not the token count\cite{27}. Stencil workloads are even more memory-bound: a 7-point FP64 stencil has an arithmetic intensity of 0.3--0.5 FLOP/byte, far below the ridge point of any current accelerator\cite{5}.

This is precisely the regime this architecture is built for. The reference chassis is a weight-stationary serving machine: every parameter and KV cache entry resides in commodity LPDDR6 pools that bypass the interposer and lithography bottlenecks constraining HBM, the router's closed-form coordinate function dispatches tokens to resident experts (Sections \hyperref[sec:2.8]{2.8} and \hyperref[sec:2.12]{2.12}), the doorbell makes each activation a sub-microsecond event (Sections \hyperref[sec:2.9]{2.9} and \hyperref[sec:3.4]{3.4}), and the spine sizing rule (Section \hyperref[sec:2.1]{2.1}) sustains $10^{8}$--$10^{9}$ fabric dispatches per second---at 270 dispatches per generated token (Section \hyperref[sec:2.6]{2.6}), a token rate of roughly $4\times10^{6}$ tokens/s at the favorable placement of Section \hyperref[sec:2.12]{2.12}, within an order of magnitude of the $\approx$10$^{7}$ tokens/s of today's largest production deployments\cite{24}. For HPC stencil workloads, the rigid X/Y grid maps directly to the motherboard topology: each grid subdomain becomes a physical node, halo exchanges are single-hop dimension-order routes, and the deterministic latency bounds enable replacing MPI with compile-time coordinate functions. A GPU serving rack uses HBM for exactly the same job---holding weights and KV cache---but at 300 W per package with capacity capped by interposer reticle limits, versus approximately 15 W per node with capacity bounded only by the number of commodity modules the chassis can hold. The roofline argument of Section \hyperref[sec:3.2]{3.2} is sharpened, not contradicted, by the serving regime: the served workload is even more memory-bound than the training-style analysis assumes.

The structural consequence is that HBM demand is coupled to exactly the workload class where this architecture substitutes commodity DRAM. Every inference token served on LPDDR6 displaces an HBM-resident weight read; because inference is the majority and fastest-growing share of AI compute\cite{21,22,23}, serving that share on capacity-optimized memory slows the industry's HBM demand growth. The trajectory is steep: inference rose from roughly a third of datacenter AI compute in 2023 to half in 2025\cite{21}, Gartner projects it overtaking training in 2026---55\% of AI-optimized IaaS spending, rising beyond 65\% by 2029\cite{22}---and McKinsey has it reaching more than half of all AI compute by 2030\cite{23}. Inference is also the workload class whose economics compound: every served token is an operating cost, token volume grows superlinearly as reasoning and agentic models multiply the work per task, and Google alone reported processing $\approx$3$\times10^{15}$ tokens per month in 2026\cite{24}. HBM's disproportionate wafer footprint---the proximate cause of the 2025--26 DRAM shortage and its price spike\cite{25}---then frees fabrication capacity back toward commodity DRAM as the mix shifts. Training retains HBM's bandwidth economics; inference---the majority---moves to capacity economics, where commodity modules with no interposer dependency and mature-node yields above 90\% dominate on both physical and economic grounds. The memory market bifurcates, and this machine is built for the inference half.

\subsection*{4.6 Security posture, design-for-test, and manufacturing readiness}\label{sec:4.6}

The security model follows from the same property that makes the fabric verifiable: state lives in exactly one place. Compute nodes hold no secrets and no mutable control state---their only configuration is the POST-loaded routing bitmap and kernel constants---so there is nothing on a node to escalate against, and the lxy\_repeater's hardwired bitmask barrier (Section \hyperref[sec:2.1]{2.1}) doubles as a hardware isolation boundary no firmware can reconfigure because no such register exists\cite{9}. The attack surface concentrates on the router chip, which the RTL addresses at four layers. First, secure boot: the 64 KB boot ROM is mask-programmed (combinational, unwritable by construction), and the firmware image it chains to is hash-verified before the CLINT timer or PNM window is released to user code. Second, privilege separation: the RV32 core implements machine-mode CSRs plus physical memory protection (PMP), and the seL4 deployment option (Section \hyperref[sec:2.17]{2.17}) replaces the monolithic kernel with a formally verified microkernel whose capability system mediates every device-frame mapping---the root task cannot touch UART or PNM registers without an explicit capability grant. Third, debug lockdown: the JTAG TAP present for bring-up fuses into a locked state at production, with re-entry requiring a signed challenge over the host PCIe link. Fourth, sideband integrity: NODE\_ERR and TOPOLOGY\_RDY travel on physically separate wires (Section \hyperref[sec:2.4]{2.4}) that no data-path fault can mask, so fault injection on the NoB cannot silence the error channel that Section \hyperref[sec:4.1]{4.1}'s recovery protocol depends on.

Design-for-test follows the same concentration principle. Both dies carry a JTAG TAP; the router SoC adds a scan-mode override on its clock-gating cells (the latch-based gate passes the test clock unconditionally when scan\_en asserts, so gated domains remain observable during scan shift), and every asynchronous reset crossing passes through a two-flop synchronizer (rst\_sync.v) so reset-domain crossings are scan-compatible rather than metastability hazards. On-chip memories---router SRAM, node elastic buffers---are reached by a bounded MBIST sequence run from ROM before POST Phase 1, so a bad SRAM word never reaches the topology inventory. The PHY blocks self-test through loopback: the PCIe endpoint's LOOPBACK parameter drives LTSSM through lane bring-up with no external partner, and the optical links' CRC-on-arrival check (Section \hyperref[sec:2.19]{2.19}) doubles as an in-field signal-integrity monitor. Manufacturing hand-off artifacts ship alongside the RTL rather than after it: SDC timing constraints for the 100 MHz reference domain with false paths enumerated for the combinational ROM/SRAM reads and async resets (pnm\_soc.sdc), a device tree source describing the full router-chip memory map for NOMMU Linux boot (pnm\_soc.dts), and the assembly schema of the companion repository that composes verified modules into board-level netlists. What remains vendor-dependent---SerDes analog characterization, package bond maps, floorplan extraction---is explicitly outside the RTL's claims, consistent with this paper's practice of labeling engineering estimates rather than claiming them away.

\section*{Conclusion}\label{sec:conclusion}

This paper has specified a distributed spatial Processing-Near-Memory architecture that bypasses the HBM capacity wall through three physical properties that no market correction can alter: commodity LPDDR6 CAMM2 modules with no interposer dependency\cite{16,17}, mature-node DUV logic with high yields and established supply chains, and deterministic wormhole routing that eliminates runtime scheduling overhead\cite{13,14}. By treating the physical motherboard as an immutable dataflow graph and exporting its topography as a bespoke ISA, the system eliminates cache coherency, operating-system overhead, and the thermal density of advanced interposer stacks. Routing is coordinate arithmetic on a single-spine tree with dimension-order on-board paths; deadlock freedom follows from four monotonically ordered virtual-channel classes with credited per-class buffers at the two state points in the design; the hardened doorbell makes activation atomic, idle-free, and loss-observable through a single ACK wire; and the isothermal parallel manifold sustains thermal viability with enforced flow balance. A RISC-V BMC/router chip at the spine root provides a programmable management plane compatible with Linux and Redox, handling topology discovery, workload dispatch, and host PCIe on a single socketed, cold-swappable die. The reference chassis---64 TB, 131 TB/s aggregate node-local bandwidth, approximately 197 TFLOPS FP64, approximately 15 W per node---demonstrates that the physical constraints of interposer packaging, EUV yield, and thermal density favour distributed commodity memory by a wide margin. At current DRAM prices, the cost advantage is approximately 20$\times$ per byte against HBM, but the structural advantage---no interposer, no EUV, no 300 W thermal hotspots---persists regardless of market cycles. The same fabric specification scales down to SODIMM-based deskside units---DDR4 or DDR5 SODIMM modules with relaxed CAS latency---and scales out to multi-chassis spatial supercomputers via a second routing level, without changing the wire protocol, the doorbell discipline, or the compiler's program schema (Section \hyperref[sec:3.7]{3.7}). Target workloads---including MoE transformer inference, FP64 stencil computation for scientific HPC, climate and seismic imaging codes, and functional data-based evaluation---sit in the bandwidth-bound regime on the precision they actually run, where mature-node silicon is already sufficient\cite{3,4,5}, and they inherit deterministic fabric replay, compile-time latency bounds, and an honest availability model with a single swappable router as the only point of programmable state.

The full repository, including Verilog HDL sources, build scripts, and this manuscript, is available at https://github.com/im-BowenGu/PNM-Arch. The hardware designs (the Verilog HDL sources) are released under the CERN Open Hardware Licence Version 2 Strongly Reciprocal (CERN-OHL-S)\cite{20}, the build and simulation code under AGPL-3.0-or-later, and this manuscript under CC BY-SA 4.0.

\section*{Acknowledgments}\label{sec:acknowledgments}

The author thanks their research mentor for guidance throughout this project, the anonymous referees for their constructive feedback, and the open-source community for developing the tools and ethos that made this work possible.


\begin{thebibliography}{30}
\bibitem{1} W. A. Wulf, S. A. McKee. Hitting the memory wall: implications of the obvious. \textit{ACM SIGARCH Computer Architecture News}. Vol. 23, pg. 20-24, 1995.
\bibitem{2} J. Choquette. Nvidia hopper h100 gpu: scaling performance. \textit{IEEE Micro}. Vol. 43, pg. 9-17, 2023.
\bibitem{3} N. Shazeer, A. Mirhoseini, K. Maziarz, A. Davis, Q. Le, G. Hinton, J. Dean. Outrageously large neural networks: the sparsely-gated mixture-of-experts layer. \textit{arXiv preprint} arXiv:1701.06538, 2017.
\bibitem{4} W. Fedus, B. Zoph, N. Shazeer. Switch transformers: scaling to trillion parameter models with simple and efficient sparsity. \textit{Journal of Machine Learning Research}. Vol. 23, pg. 1-39, 2022.
\bibitem{5} S. Williams, A. Waterman, D. Patterson. Roofline: an insightful visual performance model for multicore architectures. \textit{Communications of the ACM}. Vol. 52, pg. 65-76, 2009.
\bibitem{6} J. Backus. Can programming be liberated from the von Neumann style? a functional style and its algebra of programs. \textit{Communications of the ACM}. Vol. 21, pg. 613-641, 1978.
\bibitem{7} M. D. McIlroy, E. N. Pinson, B. A. Tague. Unix time-sharing system: forward. \textit{Bell System Technical Journal}. Vol. 57, pg. 1899-1904, 1978.
\bibitem{8} J. Liedtke. On micro-kernel construction. Proceedings of the fifteenth ACM symposium on Operating systems principles. pg. 237-250, 1995.
\bibitem{9} G. Klein, K. Elphinstone, G. Heiser, J. Andronick, D. Cock, P. Derrin, D. Elkaduwe, K. Engelhardt, R. Kolanski, M. Norrish, T. Sewell, H. Tuch, S. Winwood. seL4: formal verification of an OS kernel. Proceedings of the ACM SIGOPS 22nd symposium on Operating systems principles. pg. 207-220, 2009.
\bibitem{10} N. P. Jouppi et al. In-datacenter performance analysis of a tensor processing unit. Proceedings of the 44th annual international symposium on computer architecture. pg. 1-12, 2017.
\bibitem{11} J. B. Dennis. First version of a data flow procedure language. Programming Symposium. pg. 362-376, 1974.
\bibitem{12} I. Barron, P. Cavill, D. May, P. Wilson. Transputer does 5 or more MIPS even when not used in parallel. \textit{Electronics}. Vol. 56, pg. 109-115, 1983.
\bibitem{13} W. J. Dally, C. L. Seitz. Deadlock-free message routing in multiprocessor interconnection networks. \textit{IEEE Transactions on Computers}. Vol. C-36, pg. 547-553, 1987.
\bibitem{14} W. J. Dally, C. L. Seitz. The torus routing chip. \textit{Distributed Computing}. Vol. 1, pg. 187-196, 1986.
\bibitem{15} JEDEC Solid State Technology Association. High bandwidth memory (HBM) DRAM standard. JESD235D, 2021.
\bibitem{16} JEDEC Solid State Technology Association. Compression attached memory module (CAMM2) common standard. JESD318, 2023.
\bibitem{17} JEDEC Solid State Technology Association. Low power double data rate 6 (LPDDR6) SDRAM standard. JESD209-6, 2025.
\bibitem{18} C. Lattner et al. MLIR: scaling compiler infrastructure for domain specific computation. 2021 IEEE/ACM International Symposium on Code Generation and Optimization (CGO). pg. 2-14, 2021.
\bibitem{19} Message Passing Interface Forum. MPI: a message-passing interface standard, version 4.0. 2021.
\bibitem{20} CERN. CERN open hardware licence version 2 - strongly reciprocal (CERN-OHL-S). 2020.
\bibitem{21} Deloitte. Why AI's next phase will likely demand more computational power, not less. \textit{Deloitte Insights, TMT Predictions 2026}. 2025.
\bibitem{22} Gartner. Gartner says AI-optimized IaaS is poised to become the next growth engine for AI infrastructure. Press release, October 15, 2025.
\bibitem{23} McKinsey \& Company. The cost of compute: a \$7 trillion race to scale data centers. 2025.
\bibitem{24} R. Raghuram, S. Wang. How to win the largest market in AI. \textit{a16z}. 2026.
\bibitem{25} TrendForce. Memory spot price update: DDR5 prices up 307\% since September as module costs poised to surge. November 19, 2025.
\bibitem{26} H. Menon. Inference curves: bandwidth, not compute, is the primary serving knob. \textit{SemiAnalysis InferenceX}. 2026.
\bibitem{27} Y. Yu, et al. Efficient MoE serving in the memory-bound regime: balance activated experts, not tokens. \textit{arXiv preprint} arXiv:2512.09277, 2025.
\bibitem{28} B. Gu. Breaking the HBM wall: Verilog HDL sources, Go co-simulation harness, and manuscript repository. https://github.com/im-BowenGu/PNM-Arch, 2026.
\bibitem{29} Linux Foundation. Linux kernel RISC-V architecture documentation: NOMMU (no-MMU) support. https://www.kernel.org/doc/html/latest/arch/riscv/, 2025.
\bibitem{30} JEDEC Solid State Technology Association. Double data rate 5 (DDR5) SDRAM standard. JESD79-5C, 2024.
\end{thebibliography}

\end{document}
